\section{Technical Derivation Digest}

\label{sec:oc-core-1-3-technical-derivation-atlas}

The technical derivation corpus is preserved in the evidence package and corpus register. The public monograph prints only the reader-facing digest because a scientific manuscript should not ask the reader to traverse hundreds of generated derivation leaves before the main argument is clear.


\subsection{How the technical corpus is used}

The derivation corpus supports definitions, operator semantics, K-level examples, and falsifier routes. It is not the promoted theorem surface by itself. A promoted theorem still has to point to a proof sheet, Lean declaration when present, finite semantic witness where relevant, and a stated counterexample boundary.


\subsection{Where the full corpus lives}

The full technical derivation files remain in the release source tree and are indexed by the science monolith corpus register with hashes and inclusion roles. Reviewers who need the raw derivations should use that register and the public evidence package rather than the prose PDF.


\subsection{Technical scholarly synthesis}

The technical scholarly synthesis replaces the earlier raw atlas. It walks the same source families top-down, but each entry now states scientific role, reader use, and audit route instead of dumping unintegrated source text into the public PDF.


\subsubsection{K Levels}

\paragraph{K Levels: Klevels Master}

For K Levels: Klevels Master, the scholarly synthesis contributes K-level role, axes, thresholds, and adjacent-level witness obligations, explains why this module exists, shows how it connects to the current 1.3.3 argument, and prevents stronger inherited language from being promoted without the current proof, replay, comparator, and falsifier surfaces.


For K Levels: Klevels Master, reader use is to ask whether the level adds a non-inert observable and whether lawful demotion is available when it does not; if that source-specific question cannot be answered from the promoted theorem, proof sheet, finite semantic case, replay table, comparator row, or claim-boundary section, the source remains background support rather than public theorem evidence.


For K Levels: Klevels Master, the audit route keeps the full file `content/k\_levels/klevels\_master.tex` indexed by the science monolith corpus register and the public evidence package. The science is therefore not discarded, while the PDF avoids becoming a raw source dump. A line-level reviewer follows the register; a scientific reader follows this digest and the main chapter.


K Levels: Klevels Master reader checkpoint is satisfied only when the main manuscript tells the reader which claim class this module can support, which evidence class limits it, which prior-art or comparator boundary applies, and which failure would reopen the claim. This fourth check is deliberately positive: the module must contribute intelligibility, not merely avoid forbidden wording.


For K Levels: Klevels Master, the external wording boundary allows public language to say that the route helps organize, test, or constrain a bounded OC claim. It may not say that the route proves full-domain closure, final full-scope synthesis status, scoped unrestricted comparative claim, or a domain result absent from the current proof and evidence layer. This boundary is part of the source's scientific content, not a marketing caveat.


\paragraph{K Levels: K0}

For K Levels: K0, the scholarly synthesis contributes K-level role, axes, thresholds, and adjacent-level witness obligations, explains why this module exists, shows how it connects to the current 1.3.3 argument, and prevents stronger inherited language from being promoted without the current proof, replay, comparator, and falsifier surfaces.


For K Levels: K0, reader use is to ask whether the level adds a non-inert observable and whether lawful demotion is available when it does not; if that source-specific question cannot be answered from the promoted theorem, proof sheet, finite semantic case, replay table, comparator row, or claim-boundary section, the source remains background support rather than public theorem evidence.


For K Levels: K0, the audit route keeps the full file `content/k\_levels/k0.tex` indexed by the science monolith corpus register and the public evidence package. The science is therefore not discarded, while the PDF avoids becoming a raw source dump. A line-level reviewer follows the register; a scientific reader follows this digest and the main chapter.


K Levels: K0 reader checkpoint is satisfied only when the main manuscript tells the reader which claim class this module can support, which evidence class limits it, which prior-art or comparator boundary applies, and which failure would reopen the claim. This fourth check is deliberately positive: the module must contribute intelligibility, not merely avoid forbidden wording.


For K Levels: K0, the external wording boundary allows public language to say that the route helps organize, test, or constrain a bounded OC claim. It may not say that the route proves full-domain closure, final full-scope synthesis status, scoped unrestricted comparative claim, or a domain result absent from the current proof and evidence layer. This boundary is part of the source's scientific content, not a marketing caveat.


\paragraph{K Levels: K1}

For K Levels: K1, the scholarly synthesis contributes K-level role, axes, thresholds, and adjacent-level witness obligations, explains why this module exists, shows how it connects to the current 1.3.3 argument, and prevents stronger inherited language from being promoted without the current proof, replay, comparator, and falsifier surfaces.


For K Levels: K1, reader use is to ask whether the level adds a non-inert observable and whether lawful demotion is available when it does not; if that source-specific question cannot be answered from the promoted theorem, proof sheet, finite semantic case, replay table, comparator row, or claim-boundary section, the source remains background support rather than public theorem evidence.


For K Levels: K1, the audit route keeps the full file `content/k\_levels/k1.tex` indexed by the science monolith corpus register and the public evidence package. The science is therefore not discarded, while the PDF avoids becoming a raw source dump. A line-level reviewer follows the register; a scientific reader follows this digest and the main chapter.


K Levels: K1 reader checkpoint is satisfied only when the main manuscript tells the reader which claim class this module can support, which evidence class limits it, which prior-art or comparator boundary applies, and which failure would reopen the claim. This fourth check is deliberately positive: the module must contribute intelligibility, not merely avoid forbidden wording.


For K Levels: K1, the external wording boundary allows public language to say that the route helps organize, test, or constrain a bounded OC claim. It may not say that the route proves full-domain closure, final full-scope synthesis status, scoped unrestricted comparative claim, or a domain result absent from the current proof and evidence layer. This boundary is part of the source's scientific content, not a marketing caveat.


\paragraph{K Levels: K2}

For K Levels: K2, the scholarly synthesis contributes K-level role, axes, thresholds, and adjacent-level witness obligations, explains why this module exists, shows how it connects to the current 1.3.3 argument, and prevents stronger inherited language from being promoted without the current proof, replay, comparator, and falsifier surfaces.


For K Levels: K2, reader use is to ask whether the level adds a non-inert observable and whether lawful demotion is available when it does not; if that source-specific question cannot be answered from the promoted theorem, proof sheet, finite semantic case, replay table, comparator row, or claim-boundary section, the source remains background support rather than public theorem evidence.


For K Levels: K2, the audit route keeps the full file `content/k\_levels/k2.tex` indexed by the science monolith corpus register and the public evidence package. The science is therefore not discarded, while the PDF avoids becoming a raw source dump. A line-level reviewer follows the register; a scientific reader follows this digest and the main chapter.


K Levels: K2 reader checkpoint is satisfied only when the main manuscript tells the reader which claim class this module can support, which evidence class limits it, which prior-art or comparator boundary applies, and which failure would reopen the claim. This fourth check is deliberately positive: the module must contribute intelligibility, not merely avoid forbidden wording.


For K Levels: K2, the external wording boundary allows public language to say that the route helps organize, test, or constrain a bounded OC claim. It may not say that the route proves full-domain closure, final full-scope synthesis status, scoped unrestricted comparative claim, or a domain result absent from the current proof and evidence layer. This boundary is part of the source's scientific content, not a marketing caveat.


\paragraph{K Levels: K3}

For K Levels: K3, the scholarly synthesis contributes K-level role, axes, thresholds, and adjacent-level witness obligations, explains why this module exists, shows how it connects to the current 1.3.3 argument, and prevents stronger inherited language from being promoted without the current proof, replay, comparator, and falsifier surfaces.


For K Levels: K3, reader use is to ask whether the level adds a non-inert observable and whether lawful demotion is available when it does not; if that source-specific question cannot be answered from the promoted theorem, proof sheet, finite semantic case, replay table, comparator row, or claim-boundary section, the source remains background support rather than public theorem evidence.


For K Levels: K3, the audit route keeps the full file `content/k\_levels/k3.tex` indexed by the science monolith corpus register and the public evidence package. The science is therefore not discarded, while the PDF avoids becoming a raw source dump. A line-level reviewer follows the register; a scientific reader follows this digest and the main chapter.


K Levels: K3 reader checkpoint is satisfied only when the main manuscript tells the reader which claim class this module can support, which evidence class limits it, which prior-art or comparator boundary applies, and which failure would reopen the claim. This fourth check is deliberately positive: the module must contribute intelligibility, not merely avoid forbidden wording.


For K Levels: K3, the external wording boundary allows public language to say that the route helps organize, test, or constrain a bounded OC claim. It may not say that the route proves full-domain closure, final full-scope synthesis status, scoped unrestricted comparative claim, or a domain result absent from the current proof and evidence layer. This boundary is part of the source's scientific content, not a marketing caveat.


\paragraph{K Levels: K4}

For K Levels: K4, the scholarly synthesis contributes K-level role, axes, thresholds, and adjacent-level witness obligations, explains why this module exists, shows how it connects to the current 1.3.3 argument, and prevents stronger inherited language from being promoted without the current proof, replay, comparator, and falsifier surfaces.


For K Levels: K4, reader use is to ask whether the level adds a non-inert observable and whether lawful demotion is available when it does not; if that source-specific question cannot be answered from the promoted theorem, proof sheet, finite semantic case, replay table, comparator row, or claim-boundary section, the source remains background support rather than public theorem evidence.


For K Levels: K4, the audit route keeps the full file `content/k\_levels/k4.tex` indexed by the science monolith corpus register and the public evidence package. The science is therefore not discarded, while the PDF avoids becoming a raw source dump. A line-level reviewer follows the register; a scientific reader follows this digest and the main chapter.


K Levels: K4 reader checkpoint is satisfied only when the main manuscript tells the reader which claim class this module can support, which evidence class limits it, which prior-art or comparator boundary applies, and which failure would reopen the claim. This fourth check is deliberately positive: the module must contribute intelligibility, not merely avoid forbidden wording.


For K Levels: K4, the external wording boundary allows public language to say that the route helps organize, test, or constrain a bounded OC claim. It may not say that the route proves full-domain closure, final full-scope synthesis status, scoped unrestricted comparative claim, or a domain result absent from the current proof and evidence layer. This boundary is part of the source's scientific content, not a marketing caveat.


\paragraph{K Levels: K5}

For K Levels: K5, the scholarly synthesis contributes K-level role, axes, thresholds, and adjacent-level witness obligations, explains why this module exists, shows how it connects to the current 1.3.3 argument, and prevents stronger inherited language from being promoted without the current proof, replay, comparator, and falsifier surfaces.


For K Levels: K5, reader use is to ask whether the level adds a non-inert observable and whether lawful demotion is available when it does not; if that source-specific question cannot be answered from the promoted theorem, proof sheet, finite semantic case, replay table, comparator row, or claim-boundary section, the source remains background support rather than public theorem evidence.


For K Levels: K5, the audit route keeps the full file `content/k\_levels/k5.tex` indexed by the science monolith corpus register and the public evidence package. The science is therefore not discarded, while the PDF avoids becoming a raw source dump. A line-level reviewer follows the register; a scientific reader follows this digest and the main chapter.


K Levels: K5 reader checkpoint is satisfied only when the main manuscript tells the reader which claim class this module can support, which evidence class limits it, which prior-art or comparator boundary applies, and which failure would reopen the claim. This fourth check is deliberately positive: the module must contribute intelligibility, not merely avoid forbidden wording.


For K Levels: K5, the external wording boundary allows public language to say that the route helps organize, test, or constrain a bounded OC claim. It may not say that the route proves full-domain closure, final full-scope synthesis status, scoped unrestricted comparative claim, or a domain result absent from the current proof and evidence layer. This boundary is part of the source's scientific content, not a marketing caveat.


\paragraph{K Levels: K6}

For K Levels: K6, the scholarly synthesis contributes K-level role, axes, thresholds, and adjacent-level witness obligations, explains why this module exists, shows how it connects to the current 1.3.3 argument, and prevents stronger inherited language from being promoted without the current proof, replay, comparator, and falsifier surfaces.


For K Levels: K6, reader use is to ask whether the level adds a non-inert observable and whether lawful demotion is available when it does not; if that source-specific question cannot be answered from the promoted theorem, proof sheet, finite semantic case, replay table, comparator row, or claim-boundary section, the source remains background support rather than public theorem evidence.


For K Levels: K6, the audit route keeps the full file `content/k\_levels/k6.tex` indexed by the science monolith corpus register and the public evidence package. The science is therefore not discarded, while the PDF avoids becoming a raw source dump. A line-level reviewer follows the register; a scientific reader follows this digest and the main chapter.


K Levels: K6 reader checkpoint is satisfied only when the main manuscript tells the reader which claim class this module can support, which evidence class limits it, which prior-art or comparator boundary applies, and which failure would reopen the claim. This fourth check is deliberately positive: the module must contribute intelligibility, not merely avoid forbidden wording.


For K Levels: K6, the external wording boundary allows public language to say that the route helps organize, test, or constrain a bounded OC claim. It may not say that the route proves full-domain closure, final full-scope synthesis status, scoped unrestricted comparative claim, or a domain result absent from the current proof and evidence layer. This boundary is part of the source's scientific content, not a marketing caveat.


\paragraph{K Levels: K7}

For K Levels: K7, the scholarly synthesis contributes K-level role, axes, thresholds, and adjacent-level witness obligations, explains why this module exists, shows how it connects to the current 1.3.3 argument, and prevents stronger inherited language from being promoted without the current proof, replay, comparator, and falsifier surfaces.


For K Levels: K7, reader use is to ask whether the level adds a non-inert observable and whether lawful demotion is available when it does not; if that source-specific question cannot be answered from the promoted theorem, proof sheet, finite semantic case, replay table, comparator row, or claim-boundary section, the source remains background support rather than public theorem evidence.


For K Levels: K7, the audit route keeps the full file `content/k\_levels/k7.tex` indexed by the science monolith corpus register and the public evidence package. The science is therefore not discarded, while the PDF avoids becoming a raw source dump. A line-level reviewer follows the register; a scientific reader follows this digest and the main chapter.


K Levels: K7 reader checkpoint is satisfied only when the main manuscript tells the reader which claim class this module can support, which evidence class limits it, which prior-art or comparator boundary applies, and which failure would reopen the claim. This fourth check is deliberately positive: the module must contribute intelligibility, not merely avoid forbidden wording.


For K Levels: K7, the external wording boundary allows public language to say that the route helps organize, test, or constrain a bounded OC claim. It may not say that the route proves full-domain closure, final full-scope synthesis status, scoped unrestricted comparative claim, or a domain result absent from the current proof and evidence layer. This boundary is part of the source's scientific content, not a marketing caveat.


\paragraph{K Levels: K8}

For K Levels: K8, the scholarly synthesis contributes K-level role, axes, thresholds, and adjacent-level witness obligations, explains why this module exists, shows how it connects to the current 1.3.3 argument, and prevents stronger inherited language from being promoted without the current proof, replay, comparator, and falsifier surfaces.


For K Levels: K8, reader use is to ask whether the level adds a non-inert observable and whether lawful demotion is available when it does not; if that source-specific question cannot be answered from the promoted theorem, proof sheet, finite semantic case, replay table, comparator row, or claim-boundary section, the source remains background support rather than public theorem evidence.


For K Levels: K8, the audit route keeps the full file `content/k\_levels/k8.tex` indexed by the science monolith corpus register and the public evidence package. The science is therefore not discarded, while the PDF avoids becoming a raw source dump. A line-level reviewer follows the register; a scientific reader follows this digest and the main chapter.


K Levels: K8 reader checkpoint is satisfied only when the main manuscript tells the reader which claim class this module can support, which evidence class limits it, which prior-art or comparator boundary applies, and which failure would reopen the claim. This fourth check is deliberately positive: the module must contribute intelligibility, not merely avoid forbidden wording.


For K Levels: K8, the external wording boundary allows public language to say that the route helps organize, test, or constrain a bounded OC claim. It may not say that the route proves full-domain closure, final full-scope synthesis status, scoped unrestricted comparative claim, or a domain result absent from the current proof and evidence layer. This boundary is part of the source's scientific content, not a marketing caveat.


\paragraph{K Levels: K9}

For K Levels: K9, the scholarly synthesis contributes K-level role, axes, thresholds, and adjacent-level witness obligations, explains why this module exists, shows how it connects to the current 1.3.3 argument, and prevents stronger inherited language from being promoted without the current proof, replay, comparator, and falsifier surfaces.


For K Levels: K9, reader use is to ask whether the level adds a non-inert observable and whether lawful demotion is available when it does not; if that source-specific question cannot be answered from the promoted theorem, proof sheet, finite semantic case, replay table, comparator row, or claim-boundary section, the source remains background support rather than public theorem evidence.


For K Levels: K9, the audit route keeps the full file `content/k\_levels/k9.tex` indexed by the science monolith corpus register and the public evidence package. The science is therefore not discarded, while the PDF avoids becoming a raw source dump. A line-level reviewer follows the register; a scientific reader follows this digest and the main chapter.


K Levels: K9 reader checkpoint is satisfied only when the main manuscript tells the reader which claim class this module can support, which evidence class limits it, which prior-art or comparator boundary applies, and which failure would reopen the claim. This fourth check is deliberately positive: the module must contribute intelligibility, not merely avoid forbidden wording.


For K Levels: K9, the external wording boundary allows public language to say that the route helps organize, test, or constrain a bounded OC claim. It may not say that the route proves full-domain closure, final full-scope synthesis status, scoped unrestricted comparative claim, or a domain result absent from the current proof and evidence layer. This boundary is part of the source's scientific content, not a marketing caveat.


\paragraph{K Levels: K10}

For K Levels: K10, the scholarly synthesis contributes K-level role, axes, thresholds, and adjacent-level witness obligations, explains why this module exists, shows how it connects to the current 1.3.3 argument, and prevents stronger inherited language from being promoted without the current proof, replay, comparator, and falsifier surfaces.


For K Levels: K10, reader use is to ask whether the level adds a non-inert observable and whether lawful demotion is available when it does not; if that source-specific question cannot be answered from the promoted theorem, proof sheet, finite semantic case, replay table, comparator row, or claim-boundary section, the source remains background support rather than public theorem evidence.


For K Levels: K10, the audit route keeps the full file `content/k\_levels/k10.tex` indexed by the science monolith corpus register and the public evidence package. The science is therefore not discarded, while the PDF avoids becoming a raw source dump. A line-level reviewer follows the register; a scientific reader follows this digest and the main chapter.


K Levels: K10 reader checkpoint is satisfied only when the main manuscript tells the reader which claim class this module can support, which evidence class limits it, which prior-art or comparator boundary applies, and which failure would reopen the claim. This fourth check is deliberately positive: the module must contribute intelligibility, not merely avoid forbidden wording.


For K Levels: K10, the external wording boundary allows public language to say that the route helps organize, test, or constrain a bounded OC claim. It may not say that the route proves full-domain closure, final full-scope synthesis status, scoped unrestricted comparative claim, or a domain result absent from the current proof and evidence layer. This boundary is part of the source's scientific content, not a marketing caveat.


Synthesis checkpoint after 12 indexed modules in k\_levels. The preceding modules teach inspection logic; they do not inflate claim strength. The release remains bounded by the current proof and evidence surface.


\paragraph{K Levels: K11}

For K Levels: K11, the scholarly synthesis contributes K-level role, axes, thresholds, and adjacent-level witness obligations, explains why this module exists, shows how it connects to the current 1.3.3 argument, and prevents stronger inherited language from being promoted without the current proof, replay, comparator, and falsifier surfaces.


For K Levels: K11, reader use is to ask whether the level adds a non-inert observable and whether lawful demotion is available when it does not; if that source-specific question cannot be answered from the promoted theorem, proof sheet, finite semantic case, replay table, comparator row, or claim-boundary section, the source remains background support rather than public theorem evidence.


For K Levels: K11, the audit route keeps the full file `content/k\_levels/k11.tex` indexed by the science monolith corpus register and the public evidence package. The science is therefore not discarded, while the PDF avoids becoming a raw source dump. A line-level reviewer follows the register; a scientific reader follows this digest and the main chapter.


K Levels: K11 reader checkpoint is satisfied only when the main manuscript tells the reader which claim class this module can support, which evidence class limits it, which prior-art or comparator boundary applies, and which failure would reopen the claim. This fourth check is deliberately positive: the module must contribute intelligibility, not merely avoid forbidden wording.


For K Levels: K11, the external wording boundary allows public language to say that the route helps organize, test, or constrain a bounded OC claim. It may not say that the route proves full-domain closure, final full-scope synthesis status, scoped unrestricted comparative claim, or a domain result absent from the current proof and evidence layer. This boundary is part of the source's scientific content, not a marketing caveat.


\paragraph{K Levels: K12}

For K Levels: K12, the scholarly synthesis contributes K-level role, axes, thresholds, and adjacent-level witness obligations, explains why this module exists, shows how it connects to the current 1.3.3 argument, and prevents stronger inherited language from being promoted without the current proof, replay, comparator, and falsifier surfaces.


For K Levels: K12, reader use is to ask whether the level adds a non-inert observable and whether lawful demotion is available when it does not; if that source-specific question cannot be answered from the promoted theorem, proof sheet, finite semantic case, replay table, comparator row, or claim-boundary section, the source remains background support rather than public theorem evidence.


For K Levels: K12, the audit route keeps the full file `content/k\_levels/k12.tex` indexed by the science monolith corpus register and the public evidence package. The science is therefore not discarded, while the PDF avoids becoming a raw source dump. A line-level reviewer follows the register; a scientific reader follows this digest and the main chapter.


K Levels: K12 reader checkpoint is satisfied only when the main manuscript tells the reader which claim class this module can support, which evidence class limits it, which prior-art or comparator boundary applies, and which failure would reopen the claim. This fourth check is deliberately positive: the module must contribute intelligibility, not merely avoid forbidden wording.


For K Levels: K12, the external wording boundary allows public language to say that the route helps organize, test, or constrain a bounded OC claim. It may not say that the route proves full-domain closure, final full-scope synthesis status, scoped unrestricted comparative claim, or a domain result absent from the current proof and evidence layer. This boundary is part of the source's scientific content, not a marketing caveat.


\subsubsection{M Spaces}

\paragraph{M Spaces: Mspaces Master}

For M Spaces: Mspaces Master, the scholarly synthesis contributes embedding-space assumptions and host conditions for the level vocabulary, explains why this module exists, shows how it connects to the current 1.3.3 argument, and prevents stronger inherited language from being promoted without the current proof, replay, comparator, and falsifier surfaces.


For M Spaces: Mspaces Master, reader use is to ask whether a claimed realization has the host assumptions needed for the stated live or boundary status; if that source-specific question cannot be answered from the promoted theorem, proof sheet, finite semantic case, replay table, comparator row, or claim-boundary section, the source remains background support rather than public theorem evidence.


For M Spaces: Mspaces Master, the audit route keeps the full file `content/m\_spaces/mspaces\_master.tex` indexed by the science monolith corpus register and the public evidence package. The science is therefore not discarded, while the PDF avoids becoming a raw source dump. A line-level reviewer follows the register; a scientific reader follows this digest and the main chapter.


M Spaces: Mspaces Master reader checkpoint is satisfied only when the main manuscript tells the reader which claim class this module can support, which evidence class limits it, which prior-art or comparator boundary applies, and which failure would reopen the claim. This fourth check is deliberately positive: the module must contribute intelligibility, not merely avoid forbidden wording.


For M Spaces: Mspaces Master, the external wording boundary allows public language to say that the route helps organize, test, or constrain a bounded OC claim. It may not say that the route proves full-domain closure, final full-scope synthesis status, scoped unrestricted comparative claim, or a domain result absent from the current proof and evidence layer. This boundary is part of the source's scientific content, not a marketing caveat.


\paragraph{M Spaces: M1}

For M Spaces: M1, the scholarly synthesis contributes embedding-space assumptions and host conditions for the level vocabulary, explains why this module exists, shows how it connects to the current 1.3.3 argument, and prevents stronger inherited language from being promoted without the current proof, replay, comparator, and falsifier surfaces.


For M Spaces: M1, reader use is to ask whether a claimed realization has the host assumptions needed for the stated live or boundary status; if that source-specific question cannot be answered from the promoted theorem, proof sheet, finite semantic case, replay table, comparator row, or claim-boundary section, the source remains background support rather than public theorem evidence.


For M Spaces: M1, the audit route keeps the full file `content/m\_spaces/m1.tex` indexed by the science monolith corpus register and the public evidence package. The science is therefore not discarded, while the PDF avoids becoming a raw source dump. A line-level reviewer follows the register; a scientific reader follows this digest and the main chapter.


M Spaces: M1 reader checkpoint is satisfied only when the main manuscript tells the reader which claim class this module can support, which evidence class limits it, which prior-art or comparator boundary applies, and which failure would reopen the claim. This fourth check is deliberately positive: the module must contribute intelligibility, not merely avoid forbidden wording.


For M Spaces: M1, the external wording boundary allows public language to say that the route helps organize, test, or constrain a bounded OC claim. It may not say that the route proves full-domain closure, final full-scope synthesis status, scoped unrestricted comparative claim, or a domain result absent from the current proof and evidence layer. This boundary is part of the source's scientific content, not a marketing caveat.


\paragraph{M Spaces: M2}

For M Spaces: M2, the scholarly synthesis contributes embedding-space assumptions and host conditions for the level vocabulary, explains why this module exists, shows how it connects to the current 1.3.3 argument, and prevents stronger inherited language from being promoted without the current proof, replay, comparator, and falsifier surfaces.


For M Spaces: M2, reader use is to ask whether a claimed realization has the host assumptions needed for the stated live or boundary status; if that source-specific question cannot be answered from the promoted theorem, proof sheet, finite semantic case, replay table, comparator row, or claim-boundary section, the source remains background support rather than public theorem evidence.


For M Spaces: M2, the audit route keeps the full file `content/m\_spaces/m2.tex` indexed by the science monolith corpus register and the public evidence package. The science is therefore not discarded, while the PDF avoids becoming a raw source dump. A line-level reviewer follows the register; a scientific reader follows this digest and the main chapter.


M Spaces: M2 reader checkpoint is satisfied only when the main manuscript tells the reader which claim class this module can support, which evidence class limits it, which prior-art or comparator boundary applies, and which failure would reopen the claim. This fourth check is deliberately positive: the module must contribute intelligibility, not merely avoid forbidden wording.


For M Spaces: M2, the external wording boundary allows public language to say that the route helps organize, test, or constrain a bounded OC claim. It may not say that the route proves full-domain closure, final full-scope synthesis status, scoped unrestricted comparative claim, or a domain result absent from the current proof and evidence layer. This boundary is part of the source's scientific content, not a marketing caveat.


\paragraph{M Spaces: M3}

For M Spaces: M3, the scholarly synthesis contributes embedding-space assumptions and host conditions for the level vocabulary, explains why this module exists, shows how it connects to the current 1.3.3 argument, and prevents stronger inherited language from being promoted without the current proof, replay, comparator, and falsifier surfaces.


For M Spaces: M3, reader use is to ask whether a claimed realization has the host assumptions needed for the stated live or boundary status; if that source-specific question cannot be answered from the promoted theorem, proof sheet, finite semantic case, replay table, comparator row, or claim-boundary section, the source remains background support rather than public theorem evidence.


For M Spaces: M3, the audit route keeps the full file `content/m\_spaces/m3.tex` indexed by the science monolith corpus register and the public evidence package. The science is therefore not discarded, while the PDF avoids becoming a raw source dump. A line-level reviewer follows the register; a scientific reader follows this digest and the main chapter.


M Spaces: M3 reader checkpoint is satisfied only when the main manuscript tells the reader which claim class this module can support, which evidence class limits it, which prior-art or comparator boundary applies, and which failure would reopen the claim. This fourth check is deliberately positive: the module must contribute intelligibility, not merely avoid forbidden wording.


For M Spaces: M3, the external wording boundary allows public language to say that the route helps organize, test, or constrain a bounded OC claim. It may not say that the route proves full-domain closure, final full-scope synthesis status, scoped unrestricted comparative claim, or a domain result absent from the current proof and evidence layer. This boundary is part of the source's scientific content, not a marketing caveat.


\paragraph{M Spaces: M4}

For M Spaces: M4, the scholarly synthesis contributes embedding-space assumptions and host conditions for the level vocabulary, explains why this module exists, shows how it connects to the current 1.3.3 argument, and prevents stronger inherited language from being promoted without the current proof, replay, comparator, and falsifier surfaces.


For M Spaces: M4, reader use is to ask whether a claimed realization has the host assumptions needed for the stated live or boundary status; if that source-specific question cannot be answered from the promoted theorem, proof sheet, finite semantic case, replay table, comparator row, or claim-boundary section, the source remains background support rather than public theorem evidence.


For M Spaces: M4, the audit route keeps the full file `content/m\_spaces/m4.tex` indexed by the science monolith corpus register and the public evidence package. The science is therefore not discarded, while the PDF avoids becoming a raw source dump. A line-level reviewer follows the register; a scientific reader follows this digest and the main chapter.


M Spaces: M4 reader checkpoint is satisfied only when the main manuscript tells the reader which claim class this module can support, which evidence class limits it, which prior-art or comparator boundary applies, and which failure would reopen the claim. This fourth check is deliberately positive: the module must contribute intelligibility, not merely avoid forbidden wording.


For M Spaces: M4, the external wording boundary allows public language to say that the route helps organize, test, or constrain a bounded OC claim. It may not say that the route proves full-domain closure, final full-scope synthesis status, scoped unrestricted comparative claim, or a domain result absent from the current proof and evidence layer. This boundary is part of the source's scientific content, not a marketing caveat.


\paragraph{M Spaces: M5}

For M Spaces: M5, the scholarly synthesis contributes embedding-space assumptions and host conditions for the level vocabulary, explains why this module exists, shows how it connects to the current 1.3.3 argument, and prevents stronger inherited language from being promoted without the current proof, replay, comparator, and falsifier surfaces.


For M Spaces: M5, reader use is to ask whether a claimed realization has the host assumptions needed for the stated live or boundary status; if that source-specific question cannot be answered from the promoted theorem, proof sheet, finite semantic case, replay table, comparator row, or claim-boundary section, the source remains background support rather than public theorem evidence.


For M Spaces: M5, the audit route keeps the full file `content/m\_spaces/m5.tex` indexed by the science monolith corpus register and the public evidence package. The science is therefore not discarded, while the PDF avoids becoming a raw source dump. A line-level reviewer follows the register; a scientific reader follows this digest and the main chapter.


M Spaces: M5 reader checkpoint is satisfied only when the main manuscript tells the reader which claim class this module can support, which evidence class limits it, which prior-art or comparator boundary applies, and which failure would reopen the claim. This fourth check is deliberately positive: the module must contribute intelligibility, not merely avoid forbidden wording.


For M Spaces: M5, the external wording boundary allows public language to say that the route helps organize, test, or constrain a bounded OC claim. It may not say that the route proves full-domain closure, final full-scope synthesis status, scoped unrestricted comparative claim, or a domain result absent from the current proof and evidence layer. This boundary is part of the source's scientific content, not a marketing caveat.


\paragraph{M Spaces: M6}

For M Spaces: M6, the scholarly synthesis contributes embedding-space assumptions and host conditions for the level vocabulary, explains why this module exists, shows how it connects to the current 1.3.3 argument, and prevents stronger inherited language from being promoted without the current proof, replay, comparator, and falsifier surfaces.


For M Spaces: M6, reader use is to ask whether a claimed realization has the host assumptions needed for the stated live or boundary status; if that source-specific question cannot be answered from the promoted theorem, proof sheet, finite semantic case, replay table, comparator row, or claim-boundary section, the source remains background support rather than public theorem evidence.


For M Spaces: M6, the audit route keeps the full file `content/m\_spaces/m6.tex` indexed by the science monolith corpus register and the public evidence package. The science is therefore not discarded, while the PDF avoids becoming a raw source dump. A line-level reviewer follows the register; a scientific reader follows this digest and the main chapter.


M Spaces: M6 reader checkpoint is satisfied only when the main manuscript tells the reader which claim class this module can support, which evidence class limits it, which prior-art or comparator boundary applies, and which failure would reopen the claim. This fourth check is deliberately positive: the module must contribute intelligibility, not merely avoid forbidden wording.


For M Spaces: M6, the external wording boundary allows public language to say that the route helps organize, test, or constrain a bounded OC claim. It may not say that the route proves full-domain closure, final full-scope synthesis status, scoped unrestricted comparative claim, or a domain result absent from the current proof and evidence layer. This boundary is part of the source's scientific content, not a marketing caveat.


\paragraph{M Spaces: M7}

For M Spaces: M7, the scholarly synthesis contributes embedding-space assumptions and host conditions for the level vocabulary, explains why this module exists, shows how it connects to the current 1.3.3 argument, and prevents stronger inherited language from being promoted without the current proof, replay, comparator, and falsifier surfaces.


For M Spaces: M7, reader use is to ask whether a claimed realization has the host assumptions needed for the stated live or boundary status; if that source-specific question cannot be answered from the promoted theorem, proof sheet, finite semantic case, replay table, comparator row, or claim-boundary section, the source remains background support rather than public theorem evidence.


For M Spaces: M7, the audit route keeps the full file `content/m\_spaces/m7.tex` indexed by the science monolith corpus register and the public evidence package. The science is therefore not discarded, while the PDF avoids becoming a raw source dump. A line-level reviewer follows the register; a scientific reader follows this digest and the main chapter.


M Spaces: M7 reader checkpoint is satisfied only when the main manuscript tells the reader which claim class this module can support, which evidence class limits it, which prior-art or comparator boundary applies, and which failure would reopen the claim. This fourth check is deliberately positive: the module must contribute intelligibility, not merely avoid forbidden wording.


For M Spaces: M7, the external wording boundary allows public language to say that the route helps organize, test, or constrain a bounded OC claim. It may not say that the route proves full-domain closure, final full-scope synthesis status, scoped unrestricted comparative claim, or a domain result absent from the current proof and evidence layer. This boundary is part of the source's scientific content, not a marketing caveat.


\paragraph{M Spaces: M8}

For M Spaces: M8, the scholarly synthesis contributes embedding-space assumptions and host conditions for the level vocabulary, explains why this module exists, shows how it connects to the current 1.3.3 argument, and prevents stronger inherited language from being promoted without the current proof, replay, comparator, and falsifier surfaces.


For M Spaces: M8, reader use is to ask whether a claimed realization has the host assumptions needed for the stated live or boundary status; if that source-specific question cannot be answered from the promoted theorem, proof sheet, finite semantic case, replay table, comparator row, or claim-boundary section, the source remains background support rather than public theorem evidence.


For M Spaces: M8, the audit route keeps the full file `content/m\_spaces/m8.tex` indexed by the science monolith corpus register and the public evidence package. The science is therefore not discarded, while the PDF avoids becoming a raw source dump. A line-level reviewer follows the register; a scientific reader follows this digest and the main chapter.


M Spaces: M8 reader checkpoint is satisfied only when the main manuscript tells the reader which claim class this module can support, which evidence class limits it, which prior-art or comparator boundary applies, and which failure would reopen the claim. This fourth check is deliberately positive: the module must contribute intelligibility, not merely avoid forbidden wording.


For M Spaces: M8, the external wording boundary allows public language to say that the route helps organize, test, or constrain a bounded OC claim. It may not say that the route proves full-domain closure, final full-scope synthesis status, scoped unrestricted comparative claim, or a domain result absent from the current proof and evidence layer. This boundary is part of the source's scientific content, not a marketing caveat.


\paragraph{M Spaces: M9}

For M Spaces: M9, the scholarly synthesis contributes embedding-space assumptions and host conditions for the level vocabulary, explains why this module exists, shows how it connects to the current 1.3.3 argument, and prevents stronger inherited language from being promoted without the current proof, replay, comparator, and falsifier surfaces.


For M Spaces: M9, reader use is to ask whether a claimed realization has the host assumptions needed for the stated live or boundary status; if that source-specific question cannot be answered from the promoted theorem, proof sheet, finite semantic case, replay table, comparator row, or claim-boundary section, the source remains background support rather than public theorem evidence.


For M Spaces: M9, the audit route keeps the full file `content/m\_spaces/m9.tex` indexed by the science monolith corpus register and the public evidence package. The science is therefore not discarded, while the PDF avoids becoming a raw source dump. A line-level reviewer follows the register; a scientific reader follows this digest and the main chapter.


M Spaces: M9 reader checkpoint is satisfied only when the main manuscript tells the reader which claim class this module can support, which evidence class limits it, which prior-art or comparator boundary applies, and which failure would reopen the claim. This fourth check is deliberately positive: the module must contribute intelligibility, not merely avoid forbidden wording.


For M Spaces: M9, the external wording boundary allows public language to say that the route helps organize, test, or constrain a bounded OC claim. It may not say that the route proves full-domain closure, final full-scope synthesis status, scoped unrestricted comparative claim, or a domain result absent from the current proof and evidence layer. This boundary is part of the source's scientific content, not a marketing caveat.


Synthesis checkpoint after 24 indexed modules in m\_spaces. The preceding modules teach inspection logic; they do not inflate claim strength. The release remains bounded by the current proof and evidence surface.


\paragraph{M Spaces: M10}

For M Spaces: M10, the scholarly synthesis contributes embedding-space assumptions and host conditions for the level vocabulary, explains why this module exists, shows how it connects to the current 1.3.3 argument, and prevents stronger inherited language from being promoted without the current proof, replay, comparator, and falsifier surfaces.


For M Spaces: M10, reader use is to ask whether a claimed realization has the host assumptions needed for the stated live or boundary status; if that source-specific question cannot be answered from the promoted theorem, proof sheet, finite semantic case, replay table, comparator row, or claim-boundary section, the source remains background support rather than public theorem evidence.


For M Spaces: M10, the audit route keeps the full file `content/m\_spaces/m10.tex` indexed by the science monolith corpus register and the public evidence package. The science is therefore not discarded, while the PDF avoids becoming a raw source dump. A line-level reviewer follows the register; a scientific reader follows this digest and the main chapter.


M Spaces: M10 reader checkpoint is satisfied only when the main manuscript tells the reader which claim class this module can support, which evidence class limits it, which prior-art or comparator boundary applies, and which failure would reopen the claim. This fourth check is deliberately positive: the module must contribute intelligibility, not merely avoid forbidden wording.


For M Spaces: M10, the external wording boundary allows public language to say that the route helps organize, test, or constrain a bounded OC claim. It may not say that the route proves full-domain closure, final full-scope synthesis status, scoped unrestricted comparative claim, or a domain result absent from the current proof and evidence layer. This boundary is part of the source's scientific content, not a marketing caveat.


\paragraph{M Spaces: M11}

For M Spaces: M11, the scholarly synthesis contributes embedding-space assumptions and host conditions for the level vocabulary, explains why this module exists, shows how it connects to the current 1.3.3 argument, and prevents stronger inherited language from being promoted without the current proof, replay, comparator, and falsifier surfaces.


For M Spaces: M11, reader use is to ask whether a claimed realization has the host assumptions needed for the stated live or boundary status; if that source-specific question cannot be answered from the promoted theorem, proof sheet, finite semantic case, replay table, comparator row, or claim-boundary section, the source remains background support rather than public theorem evidence.


For M Spaces: M11, the audit route keeps the full file `content/m\_spaces/m11.tex` indexed by the science monolith corpus register and the public evidence package. The science is therefore not discarded, while the PDF avoids becoming a raw source dump. A line-level reviewer follows the register; a scientific reader follows this digest and the main chapter.


M Spaces: M11 reader checkpoint is satisfied only when the main manuscript tells the reader which claim class this module can support, which evidence class limits it, which prior-art or comparator boundary applies, and which failure would reopen the claim. This fourth check is deliberately positive: the module must contribute intelligibility, not merely avoid forbidden wording.


For M Spaces: M11, the external wording boundary allows public language to say that the route helps organize, test, or constrain a bounded OC claim. It may not say that the route proves full-domain closure, final full-scope synthesis status, scoped unrestricted comparative claim, or a domain result absent from the current proof and evidence layer. This boundary is part of the source's scientific content, not a marketing caveat.


\paragraph{M Spaces: M12}

For M Spaces: M12, the scholarly synthesis contributes embedding-space assumptions and host conditions for the level vocabulary, explains why this module exists, shows how it connects to the current 1.3.3 argument, and prevents stronger inherited language from being promoted without the current proof, replay, comparator, and falsifier surfaces.


For M Spaces: M12, reader use is to ask whether a claimed realization has the host assumptions needed for the stated live or boundary status; if that source-specific question cannot be answered from the promoted theorem, proof sheet, finite semantic case, replay table, comparator row, or claim-boundary section, the source remains background support rather than public theorem evidence.


For M Spaces: M12, the audit route keeps the full file `content/m\_spaces/m12.tex` indexed by the science monolith corpus register and the public evidence package. The science is therefore not discarded, while the PDF avoids becoming a raw source dump. A line-level reviewer follows the register; a scientific reader follows this digest and the main chapter.


M Spaces: M12 reader checkpoint is satisfied only when the main manuscript tells the reader which claim class this module can support, which evidence class limits it, which prior-art or comparator boundary applies, and which failure would reopen the claim. This fourth check is deliberately positive: the module must contribute intelligibility, not merely avoid forbidden wording.


For M Spaces: M12, the external wording boundary allows public language to say that the route helps organize, test, or constrain a bounded OC claim. It may not say that the route proves full-domain closure, final full-scope synthesis status, scoped unrestricted comparative claim, or a domain result absent from the current proof and evidence layer. This boundary is part of the source's scientific content, not a marketing caveat.


\subsubsection{Crossk}

\paragraph{Crossk: Crossk Master}

For Crossk: Crossk Master, the scholarly synthesis contributes cross-level transition semantics and reduction-failure pressure, explains why this module exists, shows how it connects to the current 1.3.3 argument, and prevents stronger inherited language from being promoted without the current proof, replay, comparator, and falsifier surfaces.


For Crossk: Crossk Master, reader use is to ask whether the proposed transition preserves a retained witness or collapses under a lower-level projection; if that source-specific question cannot be answered from the promoted theorem, proof sheet, finite semantic case, replay table, comparator row, or claim-boundary section, the source remains background support rather than public theorem evidence.


For Crossk: Crossk Master, the audit route keeps the full file `content/crossk/crossk\_master.tex` indexed by the science monolith corpus register and the public evidence package. The science is therefore not discarded, while the PDF avoids becoming a raw source dump. A line-level reviewer follows the register; a scientific reader follows this digest and the main chapter.


Crossk: Crossk Master reader checkpoint is satisfied only when the main manuscript tells the reader which claim class this module can support, which evidence class limits it, which prior-art or comparator boundary applies, and which failure would reopen the claim. This fourth check is deliberately positive: the module must contribute intelligibility, not merely avoid forbidden wording.


For Crossk: Crossk Master, the external wording boundary allows public language to say that the route helps organize, test, or constrain a bounded OC claim. It may not say that the route proves full-domain closure, final full-scope synthesis status, scoped unrestricted comparative claim, or a domain result absent from the current proof and evidence layer. This boundary is part of the source's scientific content, not a marketing caveat.


\paragraph{Crossk: Crossk Global Landscape}

For Crossk: Crossk Global Landscape, the scholarly synthesis contributes cross-level transition semantics and reduction-failure pressure, explains why this module exists, shows how it connects to the current 1.3.3 argument, and prevents stronger inherited language from being promoted without the current proof, replay, comparator, and falsifier surfaces.


For Crossk: Crossk Global Landscape, reader use is to ask whether the proposed transition preserves a retained witness or collapses under a lower-level projection; if that source-specific question cannot be answered from the promoted theorem, proof sheet, finite semantic case, replay table, comparator row, or claim-boundary section, the source remains background support rather than public theorem evidence.


For Crossk: Crossk Global Landscape, the audit route keeps the full file `content/crossk/crossk\_global\_landscape.tex` indexed by the science monolith corpus register and the public evidence package. The science is therefore not discarded, while the PDF avoids becoming a raw source dump. A line-level reviewer follows the register; a scientific reader follows this digest and the main chapter.


Crossk: Crossk Global Landscape reader checkpoint is satisfied only when the main manuscript tells the reader which claim class this module can support, which evidence class limits it, which prior-art or comparator boundary applies, and which failure would reopen the claim. This fourth check is deliberately positive: the module must contribute intelligibility, not merely avoid forbidden wording.


For Crossk: Crossk Global Landscape, the external wording boundary allows public language to say that the route helps organize, test, or constrain a bounded OC claim. It may not say that the route proves full-domain closure, final full-scope synthesis status, scoped unrestricted comparative claim, or a domain result absent from the current proof and evidence layer. This boundary is part of the source's scientific content, not a marketing caveat.


\paragraph{Crossk: Crossk K0 K1}

For Crossk: Crossk K0 K1, the scholarly synthesis contributes cross-level transition semantics and reduction-failure pressure, explains why this module exists, shows how it connects to the current 1.3.3 argument, and prevents stronger inherited language from being promoted without the current proof, replay, comparator, and falsifier surfaces.


For Crossk: Crossk K0 K1, reader use is to ask whether the proposed transition preserves a retained witness or collapses under a lower-level projection; if that source-specific question cannot be answered from the promoted theorem, proof sheet, finite semantic case, replay table, comparator row, or claim-boundary section, the source remains background support rather than public theorem evidence.


For Crossk: Crossk K0 K1, the audit route keeps the full file `content/crossk/crossk\_k0\_k1.tex` indexed by the science monolith corpus register and the public evidence package. The science is therefore not discarded, while the PDF avoids becoming a raw source dump. A line-level reviewer follows the register; a scientific reader follows this digest and the main chapter.


Crossk: Crossk K0 K1 reader checkpoint is satisfied only when the main manuscript tells the reader which claim class this module can support, which evidence class limits it, which prior-art or comparator boundary applies, and which failure would reopen the claim. This fourth check is deliberately positive: the module must contribute intelligibility, not merely avoid forbidden wording.


For Crossk: Crossk K0 K1, the external wording boundary allows public language to say that the route helps organize, test, or constrain a bounded OC claim. It may not say that the route proves full-domain closure, final full-scope synthesis status, scoped unrestricted comparative claim, or a domain result absent from the current proof and evidence layer. This boundary is part of the source's scientific content, not a marketing caveat.


\paragraph{Crossk: Crossk K1 K2}

For Crossk: Crossk K1 K2, the scholarly synthesis contributes cross-level transition semantics and reduction-failure pressure, explains why this module exists, shows how it connects to the current 1.3.3 argument, and prevents stronger inherited language from being promoted without the current proof, replay, comparator, and falsifier surfaces.


For Crossk: Crossk K1 K2, reader use is to ask whether the proposed transition preserves a retained witness or collapses under a lower-level projection; if that source-specific question cannot be answered from the promoted theorem, proof sheet, finite semantic case, replay table, comparator row, or claim-boundary section, the source remains background support rather than public theorem evidence.


For Crossk: Crossk K1 K2, the audit route keeps the full file `content/crossk/crossk\_k1\_k2.tex` indexed by the science monolith corpus register and the public evidence package. The science is therefore not discarded, while the PDF avoids becoming a raw source dump. A line-level reviewer follows the register; a scientific reader follows this digest and the main chapter.


Crossk: Crossk K1 K2 reader checkpoint is satisfied only when the main manuscript tells the reader which claim class this module can support, which evidence class limits it, which prior-art or comparator boundary applies, and which failure would reopen the claim. This fourth check is deliberately positive: the module must contribute intelligibility, not merely avoid forbidden wording.


For Crossk: Crossk K1 K2, the external wording boundary allows public language to say that the route helps organize, test, or constrain a bounded OC claim. It may not say that the route proves full-domain closure, final full-scope synthesis status, scoped unrestricted comparative claim, or a domain result absent from the current proof and evidence layer. This boundary is part of the source's scientific content, not a marketing caveat.


\paragraph{Crossk: Crossk K2 K3}

For Crossk: Crossk K2 K3, the scholarly synthesis contributes cross-level transition semantics and reduction-failure pressure, explains why this module exists, shows how it connects to the current 1.3.3 argument, and prevents stronger inherited language from being promoted without the current proof, replay, comparator, and falsifier surfaces.


For Crossk: Crossk K2 K3, reader use is to ask whether the proposed transition preserves a retained witness or collapses under a lower-level projection; if that source-specific question cannot be answered from the promoted theorem, proof sheet, finite semantic case, replay table, comparator row, or claim-boundary section, the source remains background support rather than public theorem evidence.


For Crossk: Crossk K2 K3, the audit route keeps the full file `content/crossk/crossk\_k2\_k3.tex` indexed by the science monolith corpus register and the public evidence package. The science is therefore not discarded, while the PDF avoids becoming a raw source dump. A line-level reviewer follows the register; a scientific reader follows this digest and the main chapter.


Crossk: Crossk K2 K3 reader checkpoint is satisfied only when the main manuscript tells the reader which claim class this module can support, which evidence class limits it, which prior-art or comparator boundary applies, and which failure would reopen the claim. This fourth check is deliberately positive: the module must contribute intelligibility, not merely avoid forbidden wording.


For Crossk: Crossk K2 K3, the external wording boundary allows public language to say that the route helps organize, test, or constrain a bounded OC claim. It may not say that the route proves full-domain closure, final full-scope synthesis status, scoped unrestricted comparative claim, or a domain result absent from the current proof and evidence layer. This boundary is part of the source's scientific content, not a marketing caveat.


\paragraph{Crossk: Crossk K3 K4}

For Crossk: Crossk K3 K4, the scholarly synthesis contributes cross-level transition semantics and reduction-failure pressure, explains why this module exists, shows how it connects to the current 1.3.3 argument, and prevents stronger inherited language from being promoted without the current proof, replay, comparator, and falsifier surfaces.


For Crossk: Crossk K3 K4, reader use is to ask whether the proposed transition preserves a retained witness or collapses under a lower-level projection; if that source-specific question cannot be answered from the promoted theorem, proof sheet, finite semantic case, replay table, comparator row, or claim-boundary section, the source remains background support rather than public theorem evidence.


For Crossk: Crossk K3 K4, the audit route keeps the full file `content/crossk/crossk\_k3\_k4.tex` indexed by the science monolith corpus register and the public evidence package. The science is therefore not discarded, while the PDF avoids becoming a raw source dump. A line-level reviewer follows the register; a scientific reader follows this digest and the main chapter.


Crossk: Crossk K3 K4 reader checkpoint is satisfied only when the main manuscript tells the reader which claim class this module can support, which evidence class limits it, which prior-art or comparator boundary applies, and which failure would reopen the claim. This fourth check is deliberately positive: the module must contribute intelligibility, not merely avoid forbidden wording.


For Crossk: Crossk K3 K4, the external wording boundary allows public language to say that the route helps organize, test, or constrain a bounded OC claim. It may not say that the route proves full-domain closure, final full-scope synthesis status, scoped unrestricted comparative claim, or a domain result absent from the current proof and evidence layer. This boundary is part of the source's scientific content, not a marketing caveat.


\paragraph{Crossk: Crossk K4 K5}

For Crossk: Crossk K4 K5, the scholarly synthesis contributes cross-level transition semantics and reduction-failure pressure, explains why this module exists, shows how it connects to the current 1.3.3 argument, and prevents stronger inherited language from being promoted without the current proof, replay, comparator, and falsifier surfaces.


For Crossk: Crossk K4 K5, reader use is to ask whether the proposed transition preserves a retained witness or collapses under a lower-level projection; if that source-specific question cannot be answered from the promoted theorem, proof sheet, finite semantic case, replay table, comparator row, or claim-boundary section, the source remains background support rather than public theorem evidence.


For Crossk: Crossk K4 K5, the audit route keeps the full file `content/crossk/crossk\_k4\_k5.tex` indexed by the science monolith corpus register and the public evidence package. The science is therefore not discarded, while the PDF avoids becoming a raw source dump. A line-level reviewer follows the register; a scientific reader follows this digest and the main chapter.


Crossk: Crossk K4 K5 reader checkpoint is satisfied only when the main manuscript tells the reader which claim class this module can support, which evidence class limits it, which prior-art or comparator boundary applies, and which failure would reopen the claim. This fourth check is deliberately positive: the module must contribute intelligibility, not merely avoid forbidden wording.


For Crossk: Crossk K4 K5, the external wording boundary allows public language to say that the route helps organize, test, or constrain a bounded OC claim. It may not say that the route proves full-domain closure, final full-scope synthesis status, scoped unrestricted comparative claim, or a domain result absent from the current proof and evidence layer. This boundary is part of the source's scientific content, not a marketing caveat.


\paragraph{Crossk: Crossk K5 K6}

For Crossk: Crossk K5 K6, the scholarly synthesis contributes cross-level transition semantics and reduction-failure pressure, explains why this module exists, shows how it connects to the current 1.3.3 argument, and prevents stronger inherited language from being promoted without the current proof, replay, comparator, and falsifier surfaces.


For Crossk: Crossk K5 K6, reader use is to ask whether the proposed transition preserves a retained witness or collapses under a lower-level projection; if that source-specific question cannot be answered from the promoted theorem, proof sheet, finite semantic case, replay table, comparator row, or claim-boundary section, the source remains background support rather than public theorem evidence.


For Crossk: Crossk K5 K6, the audit route keeps the full file `content/crossk/crossk\_k5\_k6.tex` indexed by the science monolith corpus register and the public evidence package. The science is therefore not discarded, while the PDF avoids becoming a raw source dump. A line-level reviewer follows the register; a scientific reader follows this digest and the main chapter.


Crossk: Crossk K5 K6 reader checkpoint is satisfied only when the main manuscript tells the reader which claim class this module can support, which evidence class limits it, which prior-art or comparator boundary applies, and which failure would reopen the claim. This fourth check is deliberately positive: the module must contribute intelligibility, not merely avoid forbidden wording.


For Crossk: Crossk K5 K6, the external wording boundary allows public language to say that the route helps organize, test, or constrain a bounded OC claim. It may not say that the route proves full-domain closure, final full-scope synthesis status, scoped unrestricted comparative claim, or a domain result absent from the current proof and evidence layer. This boundary is part of the source's scientific content, not a marketing caveat.


\paragraph{Crossk: Crossk K6 K7}

For Crossk: Crossk K6 K7, the scholarly synthesis contributes cross-level transition semantics and reduction-failure pressure, explains why this module exists, shows how it connects to the current 1.3.3 argument, and prevents stronger inherited language from being promoted without the current proof, replay, comparator, and falsifier surfaces.


For Crossk: Crossk K6 K7, reader use is to ask whether the proposed transition preserves a retained witness or collapses under a lower-level projection; if that source-specific question cannot be answered from the promoted theorem, proof sheet, finite semantic case, replay table, comparator row, or claim-boundary section, the source remains background support rather than public theorem evidence.


For Crossk: Crossk K6 K7, the audit route keeps the full file `content/crossk/crossk\_k6\_k7.tex` indexed by the science monolith corpus register and the public evidence package. The science is therefore not discarded, while the PDF avoids becoming a raw source dump. A line-level reviewer follows the register; a scientific reader follows this digest and the main chapter.


Crossk: Crossk K6 K7 reader checkpoint is satisfied only when the main manuscript tells the reader which claim class this module can support, which evidence class limits it, which prior-art or comparator boundary applies, and which failure would reopen the claim. This fourth check is deliberately positive: the module must contribute intelligibility, not merely avoid forbidden wording.


For Crossk: Crossk K6 K7, the external wording boundary allows public language to say that the route helps organize, test, or constrain a bounded OC claim. It may not say that the route proves full-domain closure, final full-scope synthesis status, scoped unrestricted comparative claim, or a domain result absent from the current proof and evidence layer. This boundary is part of the source's scientific content, not a marketing caveat.


Synthesis checkpoint after 36 indexed modules in crossk. The preceding modules teach inspection logic; they do not inflate claim strength. The release remains bounded by the current proof and evidence surface.


\paragraph{Crossk: Crossk K7 K8}

For Crossk: Crossk K7 K8, the scholarly synthesis contributes cross-level transition semantics and reduction-failure pressure, explains why this module exists, shows how it connects to the current 1.3.3 argument, and prevents stronger inherited language from being promoted without the current proof, replay, comparator, and falsifier surfaces.


For Crossk: Crossk K7 K8, reader use is to ask whether the proposed transition preserves a retained witness or collapses under a lower-level projection; if that source-specific question cannot be answered from the promoted theorem, proof sheet, finite semantic case, replay table, comparator row, or claim-boundary section, the source remains background support rather than public theorem evidence.


For Crossk: Crossk K7 K8, the audit route keeps the full file `content/crossk/crossk\_k7\_k8.tex` indexed by the science monolith corpus register and the public evidence package. The science is therefore not discarded, while the PDF avoids becoming a raw source dump. A line-level reviewer follows the register; a scientific reader follows this digest and the main chapter.


Crossk: Crossk K7 K8 reader checkpoint is satisfied only when the main manuscript tells the reader which claim class this module can support, which evidence class limits it, which prior-art or comparator boundary applies, and which failure would reopen the claim. This fourth check is deliberately positive: the module must contribute intelligibility, not merely avoid forbidden wording.


For Crossk: Crossk K7 K8, the external wording boundary allows public language to say that the route helps organize, test, or constrain a bounded OC claim. It may not say that the route proves full-domain closure, final full-scope synthesis status, scoped unrestricted comparative claim, or a domain result absent from the current proof and evidence layer. This boundary is part of the source's scientific content, not a marketing caveat.


\paragraph{Crossk: Crossk K8 K9}

For Crossk: Crossk K8 K9, the scholarly synthesis contributes cross-level transition semantics and reduction-failure pressure, explains why this module exists, shows how it connects to the current 1.3.3 argument, and prevents stronger inherited language from being promoted without the current proof, replay, comparator, and falsifier surfaces.


For Crossk: Crossk K8 K9, reader use is to ask whether the proposed transition preserves a retained witness or collapses under a lower-level projection; if that source-specific question cannot be answered from the promoted theorem, proof sheet, finite semantic case, replay table, comparator row, or claim-boundary section, the source remains background support rather than public theorem evidence.


For Crossk: Crossk K8 K9, the audit route keeps the full file `content/crossk/crossk\_k8\_k9.tex` indexed by the science monolith corpus register and the public evidence package. The science is therefore not discarded, while the PDF avoids becoming a raw source dump. A line-level reviewer follows the register; a scientific reader follows this digest and the main chapter.


Crossk: Crossk K8 K9 reader checkpoint is satisfied only when the main manuscript tells the reader which claim class this module can support, which evidence class limits it, which prior-art or comparator boundary applies, and which failure would reopen the claim. This fourth check is deliberately positive: the module must contribute intelligibility, not merely avoid forbidden wording.


For Crossk: Crossk K8 K9, the external wording boundary allows public language to say that the route helps organize, test, or constrain a bounded OC claim. It may not say that the route proves full-domain closure, final full-scope synthesis status, scoped unrestricted comparative claim, or a domain result absent from the current proof and evidence layer. This boundary is part of the source's scientific content, not a marketing caveat.


\paragraph{Crossk: Crossk K9 K10}

For Crossk: Crossk K9 K10, the scholarly synthesis contributes cross-level transition semantics and reduction-failure pressure, explains why this module exists, shows how it connects to the current 1.3.3 argument, and prevents stronger inherited language from being promoted without the current proof, replay, comparator, and falsifier surfaces.


For Crossk: Crossk K9 K10, reader use is to ask whether the proposed transition preserves a retained witness or collapses under a lower-level projection; if that source-specific question cannot be answered from the promoted theorem, proof sheet, finite semantic case, replay table, comparator row, or claim-boundary section, the source remains background support rather than public theorem evidence.


For Crossk: Crossk K9 K10, the audit route keeps the full file `content/crossk/crossk\_k9\_k10.tex` indexed by the science monolith corpus register and the public evidence package. The science is therefore not discarded, while the PDF avoids becoming a raw source dump. A line-level reviewer follows the register; a scientific reader follows this digest and the main chapter.


Crossk: Crossk K9 K10 reader checkpoint is satisfied only when the main manuscript tells the reader which claim class this module can support, which evidence class limits it, which prior-art or comparator boundary applies, and which failure would reopen the claim. This fourth check is deliberately positive: the module must contribute intelligibility, not merely avoid forbidden wording.


For Crossk: Crossk K9 K10, the external wording boundary allows public language to say that the route helps organize, test, or constrain a bounded OC claim. It may not say that the route proves full-domain closure, final full-scope synthesis status, scoped unrestricted comparative claim, or a domain result absent from the current proof and evidence layer. This boundary is part of the source's scientific content, not a marketing caveat.


\paragraph{Crossk: Crossk K10 K11}

For Crossk: Crossk K10 K11, the scholarly synthesis contributes cross-level transition semantics and reduction-failure pressure, explains why this module exists, shows how it connects to the current 1.3.3 argument, and prevents stronger inherited language from being promoted without the current proof, replay, comparator, and falsifier surfaces.


For Crossk: Crossk K10 K11, reader use is to ask whether the proposed transition preserves a retained witness or collapses under a lower-level projection; if that source-specific question cannot be answered from the promoted theorem, proof sheet, finite semantic case, replay table, comparator row, or claim-boundary section, the source remains background support rather than public theorem evidence.


For Crossk: Crossk K10 K11, the audit route keeps the full file `content/crossk/crossk\_k10\_k11.tex` indexed by the science monolith corpus register and the public evidence package. The science is therefore not discarded, while the PDF avoids becoming a raw source dump. A line-level reviewer follows the register; a scientific reader follows this digest and the main chapter.


Crossk: Crossk K10 K11 reader checkpoint is satisfied only when the main manuscript tells the reader which claim class this module can support, which evidence class limits it, which prior-art or comparator boundary applies, and which failure would reopen the claim. This fourth check is deliberately positive: the module must contribute intelligibility, not merely avoid forbidden wording.


For Crossk: Crossk K10 K11, the external wording boundary allows public language to say that the route helps organize, test, or constrain a bounded OC claim. It may not say that the route proves full-domain closure, final full-scope synthesis status, scoped unrestricted comparative claim, or a domain result absent from the current proof and evidence layer. This boundary is part of the source's scientific content, not a marketing caveat.


\paragraph{Crossk: Crossk K11 K12}

For Crossk: Crossk K11 K12, the scholarly synthesis contributes cross-level transition semantics and reduction-failure pressure, explains why this module exists, shows how it connects to the current 1.3.3 argument, and prevents stronger inherited language from being promoted without the current proof, replay, comparator, and falsifier surfaces.


For Crossk: Crossk K11 K12, reader use is to ask whether the proposed transition preserves a retained witness or collapses under a lower-level projection; if that source-specific question cannot be answered from the promoted theorem, proof sheet, finite semantic case, replay table, comparator row, or claim-boundary section, the source remains background support rather than public theorem evidence.


For Crossk: Crossk K11 K12, the audit route keeps the full file `content/crossk/crossk\_k11\_k12.tex` indexed by the science monolith corpus register and the public evidence package. The science is therefore not discarded, while the PDF avoids becoming a raw source dump. A line-level reviewer follows the register; a scientific reader follows this digest and the main chapter.


Crossk: Crossk K11 K12 reader checkpoint is satisfied only when the main manuscript tells the reader which claim class this module can support, which evidence class limits it, which prior-art or comparator boundary applies, and which failure would reopen the claim. This fourth check is deliberately positive: the module must contribute intelligibility, not merely avoid forbidden wording.


For Crossk: Crossk K11 K12, the external wording boundary allows public language to say that the route helps organize, test, or constrain a bounded OC claim. It may not say that the route proves full-domain closure, final full-scope synthesis status, scoped unrestricted comparative claim, or a domain result absent from the current proof and evidence layer. This boundary is part of the source's scientific content, not a marketing caveat.


\subsubsection{Cycles}

\paragraph{Cycles: Cycles Master}

For Cycles: Cycles Master, the scholarly synthesis contributes cycle, maintenance, and recurrence conditions, explains why this module exists, shows how it connects to the current 1.3.3 argument, and prevents stronger inherited language from being promoted without the current proof, replay, comparator, and falsifier surfaces.


For Cycles: Cycles Master, reader use is to ask whether liveness is supported by an explicit cycle mode or non-vacuous maintenance predicate; if that source-specific question cannot be answered from the promoted theorem, proof sheet, finite semantic case, replay table, comparator row, or claim-boundary section, the source remains background support rather than public theorem evidence.


For Cycles: Cycles Master, the audit route keeps the full file `content/cycles/cycles\_master.tex` indexed by the science monolith corpus register and the public evidence package. The science is therefore not discarded, while the PDF avoids becoming a raw source dump. A line-level reviewer follows the register; a scientific reader follows this digest and the main chapter.


Cycles: Cycles Master reader checkpoint is satisfied only when the main manuscript tells the reader which claim class this module can support, which evidence class limits it, which prior-art or comparator boundary applies, and which failure would reopen the claim. This fourth check is deliberately positive: the module must contribute intelligibility, not merely avoid forbidden wording.


For Cycles: Cycles Master, the external wording boundary allows public language to say that the route helps organize, test, or constrain a bounded OC claim. It may not say that the route proves full-domain closure, final full-scope synthesis status, scoped unrestricted comparative claim, or a domain result absent from the current proof and evidence layer. This boundary is part of the source's scientific content, not a marketing caveat.


\paragraph{Cycles: Cycles K0}

For Cycles: Cycles K0, the scholarly synthesis contributes cycle, maintenance, and recurrence conditions, explains why this module exists, shows how it connects to the current 1.3.3 argument, and prevents stronger inherited language from being promoted without the current proof, replay, comparator, and falsifier surfaces.


For Cycles: Cycles K0, reader use is to ask whether liveness is supported by an explicit cycle mode or non-vacuous maintenance predicate; if that source-specific question cannot be answered from the promoted theorem, proof sheet, finite semantic case, replay table, comparator row, or claim-boundary section, the source remains background support rather than public theorem evidence.


For Cycles: Cycles K0, the audit route keeps the full file `content/cycles/cycles\_k0.tex` indexed by the science monolith corpus register and the public evidence package. The science is therefore not discarded, while the PDF avoids becoming a raw source dump. A line-level reviewer follows the register; a scientific reader follows this digest and the main chapter.


Cycles: Cycles K0 reader checkpoint is satisfied only when the main manuscript tells the reader which claim class this module can support, which evidence class limits it, which prior-art or comparator boundary applies, and which failure would reopen the claim. This fourth check is deliberately positive: the module must contribute intelligibility, not merely avoid forbidden wording.


For Cycles: Cycles K0, the external wording boundary allows public language to say that the route helps organize, test, or constrain a bounded OC claim. It may not say that the route proves full-domain closure, final full-scope synthesis status, scoped unrestricted comparative claim, or a domain result absent from the current proof and evidence layer. This boundary is part of the source's scientific content, not a marketing caveat.


\paragraph{Cycles: Cycles K1}

For Cycles: Cycles K1, the scholarly synthesis contributes cycle, maintenance, and recurrence conditions, explains why this module exists, shows how it connects to the current 1.3.3 argument, and prevents stronger inherited language from being promoted without the current proof, replay, comparator, and falsifier surfaces.


For Cycles: Cycles K1, reader use is to ask whether liveness is supported by an explicit cycle mode or non-vacuous maintenance predicate; if that source-specific question cannot be answered from the promoted theorem, proof sheet, finite semantic case, replay table, comparator row, or claim-boundary section, the source remains background support rather than public theorem evidence.


For Cycles: Cycles K1, the audit route keeps the full file `content/cycles/cycles\_k1.tex` indexed by the science monolith corpus register and the public evidence package. The science is therefore not discarded, while the PDF avoids becoming a raw source dump. A line-level reviewer follows the register; a scientific reader follows this digest and the main chapter.


Cycles: Cycles K1 reader checkpoint is satisfied only when the main manuscript tells the reader which claim class this module can support, which evidence class limits it, which prior-art or comparator boundary applies, and which failure would reopen the claim. This fourth check is deliberately positive: the module must contribute intelligibility, not merely avoid forbidden wording.


For Cycles: Cycles K1, the external wording boundary allows public language to say that the route helps organize, test, or constrain a bounded OC claim. It may not say that the route proves full-domain closure, final full-scope synthesis status, scoped unrestricted comparative claim, or a domain result absent from the current proof and evidence layer. This boundary is part of the source's scientific content, not a marketing caveat.


\paragraph{Cycles: Cycles K2}

For Cycles: Cycles K2, the scholarly synthesis contributes cycle, maintenance, and recurrence conditions, explains why this module exists, shows how it connects to the current 1.3.3 argument, and prevents stronger inherited language from being promoted without the current proof, replay, comparator, and falsifier surfaces.


For Cycles: Cycles K2, reader use is to ask whether liveness is supported by an explicit cycle mode or non-vacuous maintenance predicate; if that source-specific question cannot be answered from the promoted theorem, proof sheet, finite semantic case, replay table, comparator row, or claim-boundary section, the source remains background support rather than public theorem evidence.


For Cycles: Cycles K2, the audit route keeps the full file `content/cycles/cycles\_k2.tex` indexed by the science monolith corpus register and the public evidence package. The science is therefore not discarded, while the PDF avoids becoming a raw source dump. A line-level reviewer follows the register; a scientific reader follows this digest and the main chapter.


Cycles: Cycles K2 reader checkpoint is satisfied only when the main manuscript tells the reader which claim class this module can support, which evidence class limits it, which prior-art or comparator boundary applies, and which failure would reopen the claim. This fourth check is deliberately positive: the module must contribute intelligibility, not merely avoid forbidden wording.


For Cycles: Cycles K2, the external wording boundary allows public language to say that the route helps organize, test, or constrain a bounded OC claim. It may not say that the route proves full-domain closure, final full-scope synthesis status, scoped unrestricted comparative claim, or a domain result absent from the current proof and evidence layer. This boundary is part of the source's scientific content, not a marketing caveat.


\paragraph{Cycles: Cycles K3}

For Cycles: Cycles K3, the scholarly synthesis contributes cycle, maintenance, and recurrence conditions, explains why this module exists, shows how it connects to the current 1.3.3 argument, and prevents stronger inherited language from being promoted without the current proof, replay, comparator, and falsifier surfaces.


For Cycles: Cycles K3, reader use is to ask whether liveness is supported by an explicit cycle mode or non-vacuous maintenance predicate; if that source-specific question cannot be answered from the promoted theorem, proof sheet, finite semantic case, replay table, comparator row, or claim-boundary section, the source remains background support rather than public theorem evidence.


For Cycles: Cycles K3, the audit route keeps the full file `content/cycles/cycles\_k3.tex` indexed by the science monolith corpus register and the public evidence package. The science is therefore not discarded, while the PDF avoids becoming a raw source dump. A line-level reviewer follows the register; a scientific reader follows this digest and the main chapter.


Cycles: Cycles K3 reader checkpoint is satisfied only when the main manuscript tells the reader which claim class this module can support, which evidence class limits it, which prior-art or comparator boundary applies, and which failure would reopen the claim. This fourth check is deliberately positive: the module must contribute intelligibility, not merely avoid forbidden wording.


For Cycles: Cycles K3, the external wording boundary allows public language to say that the route helps organize, test, or constrain a bounded OC claim. It may not say that the route proves full-domain closure, final full-scope synthesis status, scoped unrestricted comparative claim, or a domain result absent from the current proof and evidence layer. This boundary is part of the source's scientific content, not a marketing caveat.


\paragraph{Cycles: Cycles K4}

For Cycles: Cycles K4, the scholarly synthesis contributes cycle, maintenance, and recurrence conditions, explains why this module exists, shows how it connects to the current 1.3.3 argument, and prevents stronger inherited language from being promoted without the current proof, replay, comparator, and falsifier surfaces.


For Cycles: Cycles K4, reader use is to ask whether liveness is supported by an explicit cycle mode or non-vacuous maintenance predicate; if that source-specific question cannot be answered from the promoted theorem, proof sheet, finite semantic case, replay table, comparator row, or claim-boundary section, the source remains background support rather than public theorem evidence.


For Cycles: Cycles K4, the audit route keeps the full file `content/cycles/cycles\_k4.tex` indexed by the science monolith corpus register and the public evidence package. The science is therefore not discarded, while the PDF avoids becoming a raw source dump. A line-level reviewer follows the register; a scientific reader follows this digest and the main chapter.


Cycles: Cycles K4 reader checkpoint is satisfied only when the main manuscript tells the reader which claim class this module can support, which evidence class limits it, which prior-art or comparator boundary applies, and which failure would reopen the claim. This fourth check is deliberately positive: the module must contribute intelligibility, not merely avoid forbidden wording.


For Cycles: Cycles K4, the external wording boundary allows public language to say that the route helps organize, test, or constrain a bounded OC claim. It may not say that the route proves full-domain closure, final full-scope synthesis status, scoped unrestricted comparative claim, or a domain result absent from the current proof and evidence layer. This boundary is part of the source's scientific content, not a marketing caveat.


\paragraph{Cycles: Cycles K5}

For Cycles: Cycles K5, the scholarly synthesis contributes cycle, maintenance, and recurrence conditions, explains why this module exists, shows how it connects to the current 1.3.3 argument, and prevents stronger inherited language from being promoted without the current proof, replay, comparator, and falsifier surfaces.


For Cycles: Cycles K5, reader use is to ask whether liveness is supported by an explicit cycle mode or non-vacuous maintenance predicate; if that source-specific question cannot be answered from the promoted theorem, proof sheet, finite semantic case, replay table, comparator row, or claim-boundary section, the source remains background support rather than public theorem evidence.


For Cycles: Cycles K5, the audit route keeps the full file `content/cycles/cycles\_k5.tex` indexed by the science monolith corpus register and the public evidence package. The science is therefore not discarded, while the PDF avoids becoming a raw source dump. A line-level reviewer follows the register; a scientific reader follows this digest and the main chapter.


Cycles: Cycles K5 reader checkpoint is satisfied only when the main manuscript tells the reader which claim class this module can support, which evidence class limits it, which prior-art or comparator boundary applies, and which failure would reopen the claim. This fourth check is deliberately positive: the module must contribute intelligibility, not merely avoid forbidden wording.


For Cycles: Cycles K5, the external wording boundary allows public language to say that the route helps organize, test, or constrain a bounded OC claim. It may not say that the route proves full-domain closure, final full-scope synthesis status, scoped unrestricted comparative claim, or a domain result absent from the current proof and evidence layer. This boundary is part of the source's scientific content, not a marketing caveat.


Synthesis checkpoint after 48 indexed modules in cycles. The preceding modules teach inspection logic; they do not inflate claim strength. The release remains bounded by the current proof and evidence surface.


\paragraph{Cycles: Cycles K6}

For Cycles: Cycles K6, the scholarly synthesis contributes cycle, maintenance, and recurrence conditions, explains why this module exists, shows how it connects to the current 1.3.3 argument, and prevents stronger inherited language from being promoted without the current proof, replay, comparator, and falsifier surfaces.


For Cycles: Cycles K6, reader use is to ask whether liveness is supported by an explicit cycle mode or non-vacuous maintenance predicate; if that source-specific question cannot be answered from the promoted theorem, proof sheet, finite semantic case, replay table, comparator row, or claim-boundary section, the source remains background support rather than public theorem evidence.


For Cycles: Cycles K6, the audit route keeps the full file `content/cycles/cycles\_k6.tex` indexed by the science monolith corpus register and the public evidence package. The science is therefore not discarded, while the PDF avoids becoming a raw source dump. A line-level reviewer follows the register; a scientific reader follows this digest and the main chapter.


Cycles: Cycles K6 reader checkpoint is satisfied only when the main manuscript tells the reader which claim class this module can support, which evidence class limits it, which prior-art or comparator boundary applies, and which failure would reopen the claim. This fourth check is deliberately positive: the module must contribute intelligibility, not merely avoid forbidden wording.


For Cycles: Cycles K6, the external wording boundary allows public language to say that the route helps organize, test, or constrain a bounded OC claim. It may not say that the route proves full-domain closure, final full-scope synthesis status, scoped unrestricted comparative claim, or a domain result absent from the current proof and evidence layer. This boundary is part of the source's scientific content, not a marketing caveat.


\paragraph{Cycles: Cycles K7}

For Cycles: Cycles K7, the scholarly synthesis contributes cycle, maintenance, and recurrence conditions, explains why this module exists, shows how it connects to the current 1.3.3 argument, and prevents stronger inherited language from being promoted without the current proof, replay, comparator, and falsifier surfaces.


For Cycles: Cycles K7, reader use is to ask whether liveness is supported by an explicit cycle mode or non-vacuous maintenance predicate; if that source-specific question cannot be answered from the promoted theorem, proof sheet, finite semantic case, replay table, comparator row, or claim-boundary section, the source remains background support rather than public theorem evidence.


For Cycles: Cycles K7, the audit route keeps the full file `content/cycles/cycles\_k7.tex` indexed by the science monolith corpus register and the public evidence package. The science is therefore not discarded, while the PDF avoids becoming a raw source dump. A line-level reviewer follows the register; a scientific reader follows this digest and the main chapter.


Cycles: Cycles K7 reader checkpoint is satisfied only when the main manuscript tells the reader which claim class this module can support, which evidence class limits it, which prior-art or comparator boundary applies, and which failure would reopen the claim. This fourth check is deliberately positive: the module must contribute intelligibility, not merely avoid forbidden wording.


For Cycles: Cycles K7, the external wording boundary allows public language to say that the route helps organize, test, or constrain a bounded OC claim. It may not say that the route proves full-domain closure, final full-scope synthesis status, scoped unrestricted comparative claim, or a domain result absent from the current proof and evidence layer. This boundary is part of the source's scientific content, not a marketing caveat.


\paragraph{Cycles: Cycles K8}

For Cycles: Cycles K8, the scholarly synthesis contributes cycle, maintenance, and recurrence conditions, explains why this module exists, shows how it connects to the current 1.3.3 argument, and prevents stronger inherited language from being promoted without the current proof, replay, comparator, and falsifier surfaces.


For Cycles: Cycles K8, reader use is to ask whether liveness is supported by an explicit cycle mode or non-vacuous maintenance predicate; if that source-specific question cannot be answered from the promoted theorem, proof sheet, finite semantic case, replay table, comparator row, or claim-boundary section, the source remains background support rather than public theorem evidence.


For Cycles: Cycles K8, the audit route keeps the full file `content/cycles/cycles\_k8.tex` indexed by the science monolith corpus register and the public evidence package. The science is therefore not discarded, while the PDF avoids becoming a raw source dump. A line-level reviewer follows the register; a scientific reader follows this digest and the main chapter.


Cycles: Cycles K8 reader checkpoint is satisfied only when the main manuscript tells the reader which claim class this module can support, which evidence class limits it, which prior-art or comparator boundary applies, and which failure would reopen the claim. This fourth check is deliberately positive: the module must contribute intelligibility, not merely avoid forbidden wording.


For Cycles: Cycles K8, the external wording boundary allows public language to say that the route helps organize, test, or constrain a bounded OC claim. It may not say that the route proves full-domain closure, final full-scope synthesis status, scoped unrestricted comparative claim, or a domain result absent from the current proof and evidence layer. This boundary is part of the source's scientific content, not a marketing caveat.


\paragraph{Cycles: Cycles K9}

For Cycles: Cycles K9, the scholarly synthesis contributes cycle, maintenance, and recurrence conditions, explains why this module exists, shows how it connects to the current 1.3.3 argument, and prevents stronger inherited language from being promoted without the current proof, replay, comparator, and falsifier surfaces.


For Cycles: Cycles K9, reader use is to ask whether liveness is supported by an explicit cycle mode or non-vacuous maintenance predicate; if that source-specific question cannot be answered from the promoted theorem, proof sheet, finite semantic case, replay table, comparator row, or claim-boundary section, the source remains background support rather than public theorem evidence.


For Cycles: Cycles K9, the audit route keeps the full file `content/cycles/cycles\_k9.tex` indexed by the science monolith corpus register and the public evidence package. The science is therefore not discarded, while the PDF avoids becoming a raw source dump. A line-level reviewer follows the register; a scientific reader follows this digest and the main chapter.


Cycles: Cycles K9 reader checkpoint is satisfied only when the main manuscript tells the reader which claim class this module can support, which evidence class limits it, which prior-art or comparator boundary applies, and which failure would reopen the claim. This fourth check is deliberately positive: the module must contribute intelligibility, not merely avoid forbidden wording.


For Cycles: Cycles K9, the external wording boundary allows public language to say that the route helps organize, test, or constrain a bounded OC claim. It may not say that the route proves full-domain closure, final full-scope synthesis status, scoped unrestricted comparative claim, or a domain result absent from the current proof and evidence layer. This boundary is part of the source's scientific content, not a marketing caveat.


\paragraph{Cycles: Cycles K10}

For Cycles: Cycles K10, the scholarly synthesis contributes cycle, maintenance, and recurrence conditions, explains why this module exists, shows how it connects to the current 1.3.3 argument, and prevents stronger inherited language from being promoted without the current proof, replay, comparator, and falsifier surfaces.


For Cycles: Cycles K10, reader use is to ask whether liveness is supported by an explicit cycle mode or non-vacuous maintenance predicate; if that source-specific question cannot be answered from the promoted theorem, proof sheet, finite semantic case, replay table, comparator row, or claim-boundary section, the source remains background support rather than public theorem evidence.


For Cycles: Cycles K10, the audit route keeps the full file `content/cycles/cycles\_k10.tex` indexed by the science monolith corpus register and the public evidence package. The science is therefore not discarded, while the PDF avoids becoming a raw source dump. A line-level reviewer follows the register; a scientific reader follows this digest and the main chapter.


Cycles: Cycles K10 reader checkpoint is satisfied only when the main manuscript tells the reader which claim class this module can support, which evidence class limits it, which prior-art or comparator boundary applies, and which failure would reopen the claim. This fourth check is deliberately positive: the module must contribute intelligibility, not merely avoid forbidden wording.


For Cycles: Cycles K10, the external wording boundary allows public language to say that the route helps organize, test, or constrain a bounded OC claim. It may not say that the route proves full-domain closure, final full-scope synthesis status, scoped unrestricted comparative claim, or a domain result absent from the current proof and evidence layer. This boundary is part of the source's scientific content, not a marketing caveat.


\paragraph{Cycles: Cycles K11}

For Cycles: Cycles K11, the scholarly synthesis contributes cycle, maintenance, and recurrence conditions, explains why this module exists, shows how it connects to the current 1.3.3 argument, and prevents stronger inherited language from being promoted without the current proof, replay, comparator, and falsifier surfaces.


For Cycles: Cycles K11, reader use is to ask whether liveness is supported by an explicit cycle mode or non-vacuous maintenance predicate; if that source-specific question cannot be answered from the promoted theorem, proof sheet, finite semantic case, replay table, comparator row, or claim-boundary section, the source remains background support rather than public theorem evidence.


For Cycles: Cycles K11, the audit route keeps the full file `content/cycles/cycles\_k11.tex` indexed by the science monolith corpus register and the public evidence package. The science is therefore not discarded, while the PDF avoids becoming a raw source dump. A line-level reviewer follows the register; a scientific reader follows this digest and the main chapter.


Cycles: Cycles K11 reader checkpoint is satisfied only when the main manuscript tells the reader which claim class this module can support, which evidence class limits it, which prior-art or comparator boundary applies, and which failure would reopen the claim. This fourth check is deliberately positive: the module must contribute intelligibility, not merely avoid forbidden wording.


For Cycles: Cycles K11, the external wording boundary allows public language to say that the route helps organize, test, or constrain a bounded OC claim. It may not say that the route proves full-domain closure, final full-scope synthesis status, scoped unrestricted comparative claim, or a domain result absent from the current proof and evidence layer. This boundary is part of the source's scientific content, not a marketing caveat.


\paragraph{Cycles: Cycles K12}

For Cycles: Cycles K12, the scholarly synthesis contributes cycle, maintenance, and recurrence conditions, explains why this module exists, shows how it connects to the current 1.3.3 argument, and prevents stronger inherited language from being promoted without the current proof, replay, comparator, and falsifier surfaces.


For Cycles: Cycles K12, reader use is to ask whether liveness is supported by an explicit cycle mode or non-vacuous maintenance predicate; if that source-specific question cannot be answered from the promoted theorem, proof sheet, finite semantic case, replay table, comparator row, or claim-boundary section, the source remains background support rather than public theorem evidence.


For Cycles: Cycles K12, the audit route keeps the full file `content/cycles/cycles\_k12.tex` indexed by the science monolith corpus register and the public evidence package. The science is therefore not discarded, while the PDF avoids becoming a raw source dump. A line-level reviewer follows the register; a scientific reader follows this digest and the main chapter.


Cycles: Cycles K12 reader checkpoint is satisfied only when the main manuscript tells the reader which claim class this module can support, which evidence class limits it, which prior-art or comparator boundary applies, and which failure would reopen the claim. This fourth check is deliberately positive: the module must contribute intelligibility, not merely avoid forbidden wording.


For Cycles: Cycles K12, the external wording boundary allows public language to say that the route helps organize, test, or constrain a bounded OC claim. It may not say that the route proves full-domain closure, final full-scope synthesis status, scoped unrestricted comparative claim, or a domain result absent from the current proof and evidence layer. This boundary is part of the source's scientific content, not a marketing caveat.


\subsubsection{Jets}

\paragraph{Jets: Jets Master}

For Jets: Jets Master, the scholarly synthesis contributes flow/update information and local process structure, explains why this module exists, shows how it connects to the current 1.3.3 argument, and prevents stronger inherited language from being promoted without the current proof, replay, comparator, and falsifier surfaces.


For Jets: Jets Master, reader use is to ask whether the public claim needs smooth dynamics, discrete update semantics, or only typed transition evidence; if that source-specific question cannot be answered from the promoted theorem, proof sheet, finite semantic case, replay table, comparator row, or claim-boundary section, the source remains background support rather than public theorem evidence.


For Jets: Jets Master, the audit route keeps the full file `content/jets/jets\_master.tex` indexed by the science monolith corpus register and the public evidence package. The science is therefore not discarded, while the PDF avoids becoming a raw source dump. A line-level reviewer follows the register; a scientific reader follows this digest and the main chapter.


Jets: Jets Master reader checkpoint is satisfied only when the main manuscript tells the reader which claim class this module can support, which evidence class limits it, which prior-art or comparator boundary applies, and which failure would reopen the claim. This fourth check is deliberately positive: the module must contribute intelligibility, not merely avoid forbidden wording.


For Jets: Jets Master, the external wording boundary allows public language to say that the route helps organize, test, or constrain a bounded OC claim. It may not say that the route proves full-domain closure, final full-scope synthesis status, scoped unrestricted comparative claim, or a domain result absent from the current proof and evidence layer. This boundary is part of the source's scientific content, not a marketing caveat.


\paragraph{Jets: Jets K0}

For Jets: Jets K0, the scholarly synthesis contributes flow/update information and local process structure, explains why this module exists, shows how it connects to the current 1.3.3 argument, and prevents stronger inherited language from being promoted without the current proof, replay, comparator, and falsifier surfaces.


For Jets: Jets K0, reader use is to ask whether the public claim needs smooth dynamics, discrete update semantics, or only typed transition evidence; if that source-specific question cannot be answered from the promoted theorem, proof sheet, finite semantic case, replay table, comparator row, or claim-boundary section, the source remains background support rather than public theorem evidence.


For Jets: Jets K0, the audit route keeps the full file `content/jets/jets\_k0.tex` indexed by the science monolith corpus register and the public evidence package. The science is therefore not discarded, while the PDF avoids becoming a raw source dump. A line-level reviewer follows the register; a scientific reader follows this digest and the main chapter.


Jets: Jets K0 reader checkpoint is satisfied only when the main manuscript tells the reader which claim class this module can support, which evidence class limits it, which prior-art or comparator boundary applies, and which failure would reopen the claim. This fourth check is deliberately positive: the module must contribute intelligibility, not merely avoid forbidden wording.


For Jets: Jets K0, the external wording boundary allows public language to say that the route helps organize, test, or constrain a bounded OC claim. It may not say that the route proves full-domain closure, final full-scope synthesis status, scoped unrestricted comparative claim, or a domain result absent from the current proof and evidence layer. This boundary is part of the source's scientific content, not a marketing caveat.


\paragraph{Jets: Jets K1}

For Jets: Jets K1, the scholarly synthesis contributes flow/update information and local process structure, explains why this module exists, shows how it connects to the current 1.3.3 argument, and prevents stronger inherited language from being promoted without the current proof, replay, comparator, and falsifier surfaces.


For Jets: Jets K1, reader use is to ask whether the public claim needs smooth dynamics, discrete update semantics, or only typed transition evidence; if that source-specific question cannot be answered from the promoted theorem, proof sheet, finite semantic case, replay table, comparator row, or claim-boundary section, the source remains background support rather than public theorem evidence.


For Jets: Jets K1, the audit route keeps the full file `content/jets/jets\_k1.tex` indexed by the science monolith corpus register and the public evidence package. The science is therefore not discarded, while the PDF avoids becoming a raw source dump. A line-level reviewer follows the register; a scientific reader follows this digest and the main chapter.


Jets: Jets K1 reader checkpoint is satisfied only when the main manuscript tells the reader which claim class this module can support, which evidence class limits it, which prior-art or comparator boundary applies, and which failure would reopen the claim. This fourth check is deliberately positive: the module must contribute intelligibility, not merely avoid forbidden wording.


For Jets: Jets K1, the external wording boundary allows public language to say that the route helps organize, test, or constrain a bounded OC claim. It may not say that the route proves full-domain closure, final full-scope synthesis status, scoped unrestricted comparative claim, or a domain result absent from the current proof and evidence layer. This boundary is part of the source's scientific content, not a marketing caveat.


\paragraph{Jets: Jets K2}

For Jets: Jets K2, the scholarly synthesis contributes flow/update information and local process structure, explains why this module exists, shows how it connects to the current 1.3.3 argument, and prevents stronger inherited language from being promoted without the current proof, replay, comparator, and falsifier surfaces.


For Jets: Jets K2, reader use is to ask whether the public claim needs smooth dynamics, discrete update semantics, or only typed transition evidence; if that source-specific question cannot be answered from the promoted theorem, proof sheet, finite semantic case, replay table, comparator row, or claim-boundary section, the source remains background support rather than public theorem evidence.


For Jets: Jets K2, the audit route keeps the full file `content/jets/jets\_k2.tex` indexed by the science monolith corpus register and the public evidence package. The science is therefore not discarded, while the PDF avoids becoming a raw source dump. A line-level reviewer follows the register; a scientific reader follows this digest and the main chapter.


Jets: Jets K2 reader checkpoint is satisfied only when the main manuscript tells the reader which claim class this module can support, which evidence class limits it, which prior-art or comparator boundary applies, and which failure would reopen the claim. This fourth check is deliberately positive: the module must contribute intelligibility, not merely avoid forbidden wording.


For Jets: Jets K2, the external wording boundary allows public language to say that the route helps organize, test, or constrain a bounded OC claim. It may not say that the route proves full-domain closure, final full-scope synthesis status, scoped unrestricted comparative claim, or a domain result absent from the current proof and evidence layer. This boundary is part of the source's scientific content, not a marketing caveat.


\paragraph{Jets: Jets K3}

For Jets: Jets K3, the scholarly synthesis contributes flow/update information and local process structure, explains why this module exists, shows how it connects to the current 1.3.3 argument, and prevents stronger inherited language from being promoted without the current proof, replay, comparator, and falsifier surfaces.


For Jets: Jets K3, reader use is to ask whether the public claim needs smooth dynamics, discrete update semantics, or only typed transition evidence; if that source-specific question cannot be answered from the promoted theorem, proof sheet, finite semantic case, replay table, comparator row, or claim-boundary section, the source remains background support rather than public theorem evidence.


For Jets: Jets K3, the audit route keeps the full file `content/jets/jets\_k3.tex` indexed by the science monolith corpus register and the public evidence package. The science is therefore not discarded, while the PDF avoids becoming a raw source dump. A line-level reviewer follows the register; a scientific reader follows this digest and the main chapter.


Jets: Jets K3 reader checkpoint is satisfied only when the main manuscript tells the reader which claim class this module can support, which evidence class limits it, which prior-art or comparator boundary applies, and which failure would reopen the claim. This fourth check is deliberately positive: the module must contribute intelligibility, not merely avoid forbidden wording.


For Jets: Jets K3, the external wording boundary allows public language to say that the route helps organize, test, or constrain a bounded OC claim. It may not say that the route proves full-domain closure, final full-scope synthesis status, scoped unrestricted comparative claim, or a domain result absent from the current proof and evidence layer. This boundary is part of the source's scientific content, not a marketing caveat.


Synthesis checkpoint after 60 indexed modules in jets. The preceding modules teach inspection logic; they do not inflate claim strength. The release remains bounded by the current proof and evidence surface.


\paragraph{Jets: Jets K4}

For Jets: Jets K4, the scholarly synthesis contributes flow/update information and local process structure, explains why this module exists, shows how it connects to the current 1.3.3 argument, and prevents stronger inherited language from being promoted without the current proof, replay, comparator, and falsifier surfaces.


For Jets: Jets K4, reader use is to ask whether the public claim needs smooth dynamics, discrete update semantics, or only typed transition evidence; if that source-specific question cannot be answered from the promoted theorem, proof sheet, finite semantic case, replay table, comparator row, or claim-boundary section, the source remains background support rather than public theorem evidence.


For Jets: Jets K4, the audit route keeps the full file `content/jets/jets\_k4.tex` indexed by the science monolith corpus register and the public evidence package. The science is therefore not discarded, while the PDF avoids becoming a raw source dump. A line-level reviewer follows the register; a scientific reader follows this digest and the main chapter.


Jets: Jets K4 reader checkpoint is satisfied only when the main manuscript tells the reader which claim class this module can support, which evidence class limits it, which prior-art or comparator boundary applies, and which failure would reopen the claim. This fourth check is deliberately positive: the module must contribute intelligibility, not merely avoid forbidden wording.


For Jets: Jets K4, the external wording boundary allows public language to say that the route helps organize, test, or constrain a bounded OC claim. It may not say that the route proves full-domain closure, final full-scope synthesis status, scoped unrestricted comparative claim, or a domain result absent from the current proof and evidence layer. This boundary is part of the source's scientific content, not a marketing caveat.


\paragraph{Jets: Jets K5}

For Jets: Jets K5, the scholarly synthesis contributes flow/update information and local process structure, explains why this module exists, shows how it connects to the current 1.3.3 argument, and prevents stronger inherited language from being promoted without the current proof, replay, comparator, and falsifier surfaces.


For Jets: Jets K5, reader use is to ask whether the public claim needs smooth dynamics, discrete update semantics, or only typed transition evidence; if that source-specific question cannot be answered from the promoted theorem, proof sheet, finite semantic case, replay table, comparator row, or claim-boundary section, the source remains background support rather than public theorem evidence.


For Jets: Jets K5, the audit route keeps the full file `content/jets/jets\_k5.tex` indexed by the science monolith corpus register and the public evidence package. The science is therefore not discarded, while the PDF avoids becoming a raw source dump. A line-level reviewer follows the register; a scientific reader follows this digest and the main chapter.


Jets: Jets K5 reader checkpoint is satisfied only when the main manuscript tells the reader which claim class this module can support, which evidence class limits it, which prior-art or comparator boundary applies, and which failure would reopen the claim. This fourth check is deliberately positive: the module must contribute intelligibility, not merely avoid forbidden wording.


For Jets: Jets K5, the external wording boundary allows public language to say that the route helps organize, test, or constrain a bounded OC claim. It may not say that the route proves full-domain closure, final full-scope synthesis status, scoped unrestricted comparative claim, or a domain result absent from the current proof and evidence layer. This boundary is part of the source's scientific content, not a marketing caveat.


\paragraph{Jets: Jets K6}

For Jets: Jets K6, the scholarly synthesis contributes flow/update information and local process structure, explains why this module exists, shows how it connects to the current 1.3.3 argument, and prevents stronger inherited language from being promoted without the current proof, replay, comparator, and falsifier surfaces.


For Jets: Jets K6, reader use is to ask whether the public claim needs smooth dynamics, discrete update semantics, or only typed transition evidence; if that source-specific question cannot be answered from the promoted theorem, proof sheet, finite semantic case, replay table, comparator row, or claim-boundary section, the source remains background support rather than public theorem evidence.


For Jets: Jets K6, the audit route keeps the full file `content/jets/jets\_k6.tex` indexed by the science monolith corpus register and the public evidence package. The science is therefore not discarded, while the PDF avoids becoming a raw source dump. A line-level reviewer follows the register; a scientific reader follows this digest and the main chapter.


Jets: Jets K6 reader checkpoint is satisfied only when the main manuscript tells the reader which claim class this module can support, which evidence class limits it, which prior-art or comparator boundary applies, and which failure would reopen the claim. This fourth check is deliberately positive: the module must contribute intelligibility, not merely avoid forbidden wording.


For Jets: Jets K6, the external wording boundary allows public language to say that the route helps organize, test, or constrain a bounded OC claim. It may not say that the route proves full-domain closure, final full-scope synthesis status, scoped unrestricted comparative claim, or a domain result absent from the current proof and evidence layer. This boundary is part of the source's scientific content, not a marketing caveat.


\paragraph{Jets: Jets K7}

For Jets: Jets K7, the scholarly synthesis contributes flow/update information and local process structure, explains why this module exists, shows how it connects to the current 1.3.3 argument, and prevents stronger inherited language from being promoted without the current proof, replay, comparator, and falsifier surfaces.


For Jets: Jets K7, reader use is to ask whether the public claim needs smooth dynamics, discrete update semantics, or only typed transition evidence; if that source-specific question cannot be answered from the promoted theorem, proof sheet, finite semantic case, replay table, comparator row, or claim-boundary section, the source remains background support rather than public theorem evidence.


For Jets: Jets K7, the audit route keeps the full file `content/jets/jets\_k7.tex` indexed by the science monolith corpus register and the public evidence package. The science is therefore not discarded, while the PDF avoids becoming a raw source dump. A line-level reviewer follows the register; a scientific reader follows this digest and the main chapter.


Jets: Jets K7 reader checkpoint is satisfied only when the main manuscript tells the reader which claim class this module can support, which evidence class limits it, which prior-art or comparator boundary applies, and which failure would reopen the claim. This fourth check is deliberately positive: the module must contribute intelligibility, not merely avoid forbidden wording.


For Jets: Jets K7, the external wording boundary allows public language to say that the route helps organize, test, or constrain a bounded OC claim. It may not say that the route proves full-domain closure, final full-scope synthesis status, scoped unrestricted comparative claim, or a domain result absent from the current proof and evidence layer. This boundary is part of the source's scientific content, not a marketing caveat.


\paragraph{Jets: Jets K8}

For Jets: Jets K8, the scholarly synthesis contributes flow/update information and local process structure, explains why this module exists, shows how it connects to the current 1.3.3 argument, and prevents stronger inherited language from being promoted without the current proof, replay, comparator, and falsifier surfaces.


For Jets: Jets K8, reader use is to ask whether the public claim needs smooth dynamics, discrete update semantics, or only typed transition evidence; if that source-specific question cannot be answered from the promoted theorem, proof sheet, finite semantic case, replay table, comparator row, or claim-boundary section, the source remains background support rather than public theorem evidence.


For Jets: Jets K8, the audit route keeps the full file `content/jets/jets\_k8.tex` indexed by the science monolith corpus register and the public evidence package. The science is therefore not discarded, while the PDF avoids becoming a raw source dump. A line-level reviewer follows the register; a scientific reader follows this digest and the main chapter.


Jets: Jets K8 reader checkpoint is satisfied only when the main manuscript tells the reader which claim class this module can support, which evidence class limits it, which prior-art or comparator boundary applies, and which failure would reopen the claim. This fourth check is deliberately positive: the module must contribute intelligibility, not merely avoid forbidden wording.


For Jets: Jets K8, the external wording boundary allows public language to say that the route helps organize, test, or constrain a bounded OC claim. It may not say that the route proves full-domain closure, final full-scope synthesis status, scoped unrestricted comparative claim, or a domain result absent from the current proof and evidence layer. This boundary is part of the source's scientific content, not a marketing caveat.


\paragraph{Jets: Jets K9}

For Jets: Jets K9, the scholarly synthesis contributes flow/update information and local process structure, explains why this module exists, shows how it connects to the current 1.3.3 argument, and prevents stronger inherited language from being promoted without the current proof, replay, comparator, and falsifier surfaces.


For Jets: Jets K9, reader use is to ask whether the public claim needs smooth dynamics, discrete update semantics, or only typed transition evidence; if that source-specific question cannot be answered from the promoted theorem, proof sheet, finite semantic case, replay table, comparator row, or claim-boundary section, the source remains background support rather than public theorem evidence.


For Jets: Jets K9, the audit route keeps the full file `content/jets/jets\_k9.tex` indexed by the science monolith corpus register and the public evidence package. The science is therefore not discarded, while the PDF avoids becoming a raw source dump. A line-level reviewer follows the register; a scientific reader follows this digest and the main chapter.


Jets: Jets K9 reader checkpoint is satisfied only when the main manuscript tells the reader which claim class this module can support, which evidence class limits it, which prior-art or comparator boundary applies, and which failure would reopen the claim. This fourth check is deliberately positive: the module must contribute intelligibility, not merely avoid forbidden wording.


For Jets: Jets K9, the external wording boundary allows public language to say that the route helps organize, test, or constrain a bounded OC claim. It may not say that the route proves full-domain closure, final full-scope synthesis status, scoped unrestricted comparative claim, or a domain result absent from the current proof and evidence layer. This boundary is part of the source's scientific content, not a marketing caveat.


\paragraph{Jets: Jets K10}

For Jets: Jets K10, the scholarly synthesis contributes flow/update information and local process structure, explains why this module exists, shows how it connects to the current 1.3.3 argument, and prevents stronger inherited language from being promoted without the current proof, replay, comparator, and falsifier surfaces.


For Jets: Jets K10, reader use is to ask whether the public claim needs smooth dynamics, discrete update semantics, or only typed transition evidence; if that source-specific question cannot be answered from the promoted theorem, proof sheet, finite semantic case, replay table, comparator row, or claim-boundary section, the source remains background support rather than public theorem evidence.


For Jets: Jets K10, the audit route keeps the full file `content/jets/jets\_k10.tex` indexed by the science monolith corpus register and the public evidence package. The science is therefore not discarded, while the PDF avoids becoming a raw source dump. A line-level reviewer follows the register; a scientific reader follows this digest and the main chapter.


Jets: Jets K10 reader checkpoint is satisfied only when the main manuscript tells the reader which claim class this module can support, which evidence class limits it, which prior-art or comparator boundary applies, and which failure would reopen the claim. This fourth check is deliberately positive: the module must contribute intelligibility, not merely avoid forbidden wording.


For Jets: Jets K10, the external wording boundary allows public language to say that the route helps organize, test, or constrain a bounded OC claim. It may not say that the route proves full-domain closure, final full-scope synthesis status, scoped unrestricted comparative claim, or a domain result absent from the current proof and evidence layer. This boundary is part of the source's scientific content, not a marketing caveat.


\paragraph{Jets: Jets K11}

For Jets: Jets K11, the scholarly synthesis contributes flow/update information and local process structure, explains why this module exists, shows how it connects to the current 1.3.3 argument, and prevents stronger inherited language from being promoted without the current proof, replay, comparator, and falsifier surfaces.


For Jets: Jets K11, reader use is to ask whether the public claim needs smooth dynamics, discrete update semantics, or only typed transition evidence; if that source-specific question cannot be answered from the promoted theorem, proof sheet, finite semantic case, replay table, comparator row, or claim-boundary section, the source remains background support rather than public theorem evidence.


For Jets: Jets K11, the audit route keeps the full file `content/jets/jets\_k11.tex` indexed by the science monolith corpus register and the public evidence package. The science is therefore not discarded, while the PDF avoids becoming a raw source dump. A line-level reviewer follows the register; a scientific reader follows this digest and the main chapter.


Jets: Jets K11 reader checkpoint is satisfied only when the main manuscript tells the reader which claim class this module can support, which evidence class limits it, which prior-art or comparator boundary applies, and which failure would reopen the claim. This fourth check is deliberately positive: the module must contribute intelligibility, not merely avoid forbidden wording.


For Jets: Jets K11, the external wording boundary allows public language to say that the route helps organize, test, or constrain a bounded OC claim. It may not say that the route proves full-domain closure, final full-scope synthesis status, scoped unrestricted comparative claim, or a domain result absent from the current proof and evidence layer. This boundary is part of the source's scientific content, not a marketing caveat.


\paragraph{Jets: Jets K12}

For Jets: Jets K12, the scholarly synthesis contributes flow/update information and local process structure, explains why this module exists, shows how it connects to the current 1.3.3 argument, and prevents stronger inherited language from being promoted without the current proof, replay, comparator, and falsifier surfaces.


For Jets: Jets K12, reader use is to ask whether the public claim needs smooth dynamics, discrete update semantics, or only typed transition evidence; if that source-specific question cannot be answered from the promoted theorem, proof sheet, finite semantic case, replay table, comparator row, or claim-boundary section, the source remains background support rather than public theorem evidence.


For Jets: Jets K12, the audit route keeps the full file `content/jets/jets\_k12.tex` indexed by the science monolith corpus register and the public evidence package. The science is therefore not discarded, while the PDF avoids becoming a raw source dump. A line-level reviewer follows the register; a scientific reader follows this digest and the main chapter.


Jets: Jets K12 reader checkpoint is satisfied only when the main manuscript tells the reader which claim class this module can support, which evidence class limits it, which prior-art or comparator boundary applies, and which failure would reopen the claim. This fourth check is deliberately positive: the module must contribute intelligibility, not merely avoid forbidden wording.


For Jets: Jets K12, the external wording boundary allows public language to say that the route helps organize, test, or constrain a bounded OC claim. It may not say that the route proves full-domain closure, final full-scope synthesis status, scoped unrestricted comparative claim, or a domain result absent from the current proof and evidence layer. This boundary is part of the source's scientific content, not a marketing caveat.


\subsubsection{Processes}

\paragraph{Processes: Processes Master}

For Processes: Processes Master, the scholarly synthesis contributes process schema, ordering, and change conditions, explains why this module exists, shows how it connects to the current 1.3.3 argument, and prevents stronger inherited language from being promoted without the current proof, replay, comparator, and falsifier surfaces.


For Processes: Processes Master, reader use is to ask whether the process claim names the state, transition, boundary, and endpoint condition that make it testable; if that source-specific question cannot be answered from the promoted theorem, proof sheet, finite semantic case, replay table, comparator row, or claim-boundary section, the source remains background support rather than public theorem evidence.


For Processes: Processes Master, the audit route keeps the full file `content/processes/processes\_master.tex` indexed by the science monolith corpus register and the public evidence package. The science is therefore not discarded, while the PDF avoids becoming a raw source dump. A line-level reviewer follows the register; a scientific reader follows this digest and the main chapter.


Processes: Processes Master reader checkpoint is satisfied only when the main manuscript tells the reader which claim class this module can support, which evidence class limits it, which prior-art or comparator boundary applies, and which failure would reopen the claim. This fourth check is deliberately positive: the module must contribute intelligibility, not merely avoid forbidden wording.


For Processes: Processes Master, the external wording boundary allows public language to say that the route helps organize, test, or constrain a bounded OC claim. It may not say that the route proves full-domain closure, final full-scope synthesis status, scoped unrestricted comparative claim, or a domain result absent from the current proof and evidence layer. This boundary is part of the source's scientific content, not a marketing caveat.


\paragraph{Processes: Processes K0}

For Processes: Processes K0, the scholarly synthesis contributes process schema, ordering, and change conditions, explains why this module exists, shows how it connects to the current 1.3.3 argument, and prevents stronger inherited language from being promoted without the current proof, replay, comparator, and falsifier surfaces.


For Processes: Processes K0, reader use is to ask whether the process claim names the state, transition, boundary, and endpoint condition that make it testable; if that source-specific question cannot be answered from the promoted theorem, proof sheet, finite semantic case, replay table, comparator row, or claim-boundary section, the source remains background support rather than public theorem evidence.


For Processes: Processes K0, the audit route keeps the full file `content/processes/processes\_k0.tex` indexed by the science monolith corpus register and the public evidence package. The science is therefore not discarded, while the PDF avoids becoming a raw source dump. A line-level reviewer follows the register; a scientific reader follows this digest and the main chapter.


Processes: Processes K0 reader checkpoint is satisfied only when the main manuscript tells the reader which claim class this module can support, which evidence class limits it, which prior-art or comparator boundary applies, and which failure would reopen the claim. This fourth check is deliberately positive: the module must contribute intelligibility, not merely avoid forbidden wording.


For Processes: Processes K0, the external wording boundary allows public language to say that the route helps organize, test, or constrain a bounded OC claim. It may not say that the route proves full-domain closure, final full-scope synthesis status, scoped unrestricted comparative claim, or a domain result absent from the current proof and evidence layer. This boundary is part of the source's scientific content, not a marketing caveat.


\paragraph{Processes: Processes K1}

For Processes: Processes K1, the scholarly synthesis contributes process schema, ordering, and change conditions, explains why this module exists, shows how it connects to the current 1.3.3 argument, and prevents stronger inherited language from being promoted without the current proof, replay, comparator, and falsifier surfaces.


For Processes: Processes K1, reader use is to ask whether the process claim names the state, transition, boundary, and endpoint condition that make it testable; if that source-specific question cannot be answered from the promoted theorem, proof sheet, finite semantic case, replay table, comparator row, or claim-boundary section, the source remains background support rather than public theorem evidence.


For Processes: Processes K1, the audit route keeps the full file `content/processes/processes\_k1.tex` indexed by the science monolith corpus register and the public evidence package. The science is therefore not discarded, while the PDF avoids becoming a raw source dump. A line-level reviewer follows the register; a scientific reader follows this digest and the main chapter.


Processes: Processes K1 reader checkpoint is satisfied only when the main manuscript tells the reader which claim class this module can support, which evidence class limits it, which prior-art or comparator boundary applies, and which failure would reopen the claim. This fourth check is deliberately positive: the module must contribute intelligibility, not merely avoid forbidden wording.


For Processes: Processes K1, the external wording boundary allows public language to say that the route helps organize, test, or constrain a bounded OC claim. It may not say that the route proves full-domain closure, final full-scope synthesis status, scoped unrestricted comparative claim, or a domain result absent from the current proof and evidence layer. This boundary is part of the source's scientific content, not a marketing caveat.


Synthesis checkpoint after 72 indexed modules in processes. The preceding modules teach inspection logic; they do not inflate claim strength. The release remains bounded by the current proof and evidence surface.


\paragraph{Processes: Processes K2}

For Processes: Processes K2, the scholarly synthesis contributes process schema, ordering, and change conditions, explains why this module exists, shows how it connects to the current 1.3.3 argument, and prevents stronger inherited language from being promoted without the current proof, replay, comparator, and falsifier surfaces.


For Processes: Processes K2, reader use is to ask whether the process claim names the state, transition, boundary, and endpoint condition that make it testable; if that source-specific question cannot be answered from the promoted theorem, proof sheet, finite semantic case, replay table, comparator row, or claim-boundary section, the source remains background support rather than public theorem evidence.


For Processes: Processes K2, the audit route keeps the full file `content/processes/processes\_k2.tex` indexed by the science monolith corpus register and the public evidence package. The science is therefore not discarded, while the PDF avoids becoming a raw source dump. A line-level reviewer follows the register; a scientific reader follows this digest and the main chapter.


Processes: Processes K2 reader checkpoint is satisfied only when the main manuscript tells the reader which claim class this module can support, which evidence class limits it, which prior-art or comparator boundary applies, and which failure would reopen the claim. This fourth check is deliberately positive: the module must contribute intelligibility, not merely avoid forbidden wording.


For Processes: Processes K2, the external wording boundary allows public language to say that the route helps organize, test, or constrain a bounded OC claim. It may not say that the route proves full-domain closure, final full-scope synthesis status, scoped unrestricted comparative claim, or a domain result absent from the current proof and evidence layer. This boundary is part of the source's scientific content, not a marketing caveat.


\paragraph{Processes: Processes K3}

For Processes: Processes K3, the scholarly synthesis contributes process schema, ordering, and change conditions, explains why this module exists, shows how it connects to the current 1.3.3 argument, and prevents stronger inherited language from being promoted without the current proof, replay, comparator, and falsifier surfaces.


For Processes: Processes K3, reader use is to ask whether the process claim names the state, transition, boundary, and endpoint condition that make it testable; if that source-specific question cannot be answered from the promoted theorem, proof sheet, finite semantic case, replay table, comparator row, or claim-boundary section, the source remains background support rather than public theorem evidence.


For Processes: Processes K3, the audit route keeps the full file `content/processes/processes\_k3.tex` indexed by the science monolith corpus register and the public evidence package. The science is therefore not discarded, while the PDF avoids becoming a raw source dump. A line-level reviewer follows the register; a scientific reader follows this digest and the main chapter.


Processes: Processes K3 reader checkpoint is satisfied only when the main manuscript tells the reader which claim class this module can support, which evidence class limits it, which prior-art or comparator boundary applies, and which failure would reopen the claim. This fourth check is deliberately positive: the module must contribute intelligibility, not merely avoid forbidden wording.


For Processes: Processes K3, the external wording boundary allows public language to say that the route helps organize, test, or constrain a bounded OC claim. It may not say that the route proves full-domain closure, final full-scope synthesis status, scoped unrestricted comparative claim, or a domain result absent from the current proof and evidence layer. This boundary is part of the source's scientific content, not a marketing caveat.


\paragraph{Processes: Processes K4}

For Processes: Processes K4, the scholarly synthesis contributes process schema, ordering, and change conditions, explains why this module exists, shows how it connects to the current 1.3.3 argument, and prevents stronger inherited language from being promoted without the current proof, replay, comparator, and falsifier surfaces.


For Processes: Processes K4, reader use is to ask whether the process claim names the state, transition, boundary, and endpoint condition that make it testable; if that source-specific question cannot be answered from the promoted theorem, proof sheet, finite semantic case, replay table, comparator row, or claim-boundary section, the source remains background support rather than public theorem evidence.


For Processes: Processes K4, the audit route keeps the full file `content/processes/processes\_k4.tex` indexed by the science monolith corpus register and the public evidence package. The science is therefore not discarded, while the PDF avoids becoming a raw source dump. A line-level reviewer follows the register; a scientific reader follows this digest and the main chapter.


Processes: Processes K4 reader checkpoint is satisfied only when the main manuscript tells the reader which claim class this module can support, which evidence class limits it, which prior-art or comparator boundary applies, and which failure would reopen the claim. This fourth check is deliberately positive: the module must contribute intelligibility, not merely avoid forbidden wording.


For Processes: Processes K4, the external wording boundary allows public language to say that the route helps organize, test, or constrain a bounded OC claim. It may not say that the route proves full-domain closure, final full-scope synthesis status, scoped unrestricted comparative claim, or a domain result absent from the current proof and evidence layer. This boundary is part of the source's scientific content, not a marketing caveat.


\paragraph{Processes: Processes K5}

For Processes: Processes K5, the scholarly synthesis contributes process schema, ordering, and change conditions, explains why this module exists, shows how it connects to the current 1.3.3 argument, and prevents stronger inherited language from being promoted without the current proof, replay, comparator, and falsifier surfaces.


For Processes: Processes K5, reader use is to ask whether the process claim names the state, transition, boundary, and endpoint condition that make it testable; if that source-specific question cannot be answered from the promoted theorem, proof sheet, finite semantic case, replay table, comparator row, or claim-boundary section, the source remains background support rather than public theorem evidence.


For Processes: Processes K5, the audit route keeps the full file `content/processes/processes\_k5.tex` indexed by the science monolith corpus register and the public evidence package. The science is therefore not discarded, while the PDF avoids becoming a raw source dump. A line-level reviewer follows the register; a scientific reader follows this digest and the main chapter.


Processes: Processes K5 reader checkpoint is satisfied only when the main manuscript tells the reader which claim class this module can support, which evidence class limits it, which prior-art or comparator boundary applies, and which failure would reopen the claim. This fourth check is deliberately positive: the module must contribute intelligibility, not merely avoid forbidden wording.


For Processes: Processes K5, the external wording boundary allows public language to say that the route helps organize, test, or constrain a bounded OC claim. It may not say that the route proves full-domain closure, final full-scope synthesis status, scoped unrestricted comparative claim, or a domain result absent from the current proof and evidence layer. This boundary is part of the source's scientific content, not a marketing caveat.


\paragraph{Processes: Processes K6}

For Processes: Processes K6, the scholarly synthesis contributes process schema, ordering, and change conditions, explains why this module exists, shows how it connects to the current 1.3.3 argument, and prevents stronger inherited language from being promoted without the current proof, replay, comparator, and falsifier surfaces.


For Processes: Processes K6, reader use is to ask whether the process claim names the state, transition, boundary, and endpoint condition that make it testable; if that source-specific question cannot be answered from the promoted theorem, proof sheet, finite semantic case, replay table, comparator row, or claim-boundary section, the source remains background support rather than public theorem evidence.


For Processes: Processes K6, the audit route keeps the full file `content/processes/processes\_k6.tex` indexed by the science monolith corpus register and the public evidence package. The science is therefore not discarded, while the PDF avoids becoming a raw source dump. A line-level reviewer follows the register; a scientific reader follows this digest and the main chapter.


Processes: Processes K6 reader checkpoint is satisfied only when the main manuscript tells the reader which claim class this module can support, which evidence class limits it, which prior-art or comparator boundary applies, and which failure would reopen the claim. This fourth check is deliberately positive: the module must contribute intelligibility, not merely avoid forbidden wording.


For Processes: Processes K6, the external wording boundary allows public language to say that the route helps organize, test, or constrain a bounded OC claim. It may not say that the route proves full-domain closure, final full-scope synthesis status, scoped unrestricted comparative claim, or a domain result absent from the current proof and evidence layer. This boundary is part of the source's scientific content, not a marketing caveat.


\paragraph{Processes: Processes K7}

For Processes: Processes K7, the scholarly synthesis contributes process schema, ordering, and change conditions, explains why this module exists, shows how it connects to the current 1.3.3 argument, and prevents stronger inherited language from being promoted without the current proof, replay, comparator, and falsifier surfaces.


For Processes: Processes K7, reader use is to ask whether the process claim names the state, transition, boundary, and endpoint condition that make it testable; if that source-specific question cannot be answered from the promoted theorem, proof sheet, finite semantic case, replay table, comparator row, or claim-boundary section, the source remains background support rather than public theorem evidence.


For Processes: Processes K7, the audit route keeps the full file `content/processes/processes\_k7.tex` indexed by the science monolith corpus register and the public evidence package. The science is therefore not discarded, while the PDF avoids becoming a raw source dump. A line-level reviewer follows the register; a scientific reader follows this digest and the main chapter.


Processes: Processes K7 reader checkpoint is satisfied only when the main manuscript tells the reader which claim class this module can support, which evidence class limits it, which prior-art or comparator boundary applies, and which failure would reopen the claim. This fourth check is deliberately positive: the module must contribute intelligibility, not merely avoid forbidden wording.


For Processes: Processes K7, the external wording boundary allows public language to say that the route helps organize, test, or constrain a bounded OC claim. It may not say that the route proves full-domain closure, final full-scope synthesis status, scoped unrestricted comparative claim, or a domain result absent from the current proof and evidence layer. This boundary is part of the source's scientific content, not a marketing caveat.


\paragraph{Processes: Processes K8}

For Processes: Processes K8, the scholarly synthesis contributes process schema, ordering, and change conditions, explains why this module exists, shows how it connects to the current 1.3.3 argument, and prevents stronger inherited language from being promoted without the current proof, replay, comparator, and falsifier surfaces.


For Processes: Processes K8, reader use is to ask whether the process claim names the state, transition, boundary, and endpoint condition that make it testable; if that source-specific question cannot be answered from the promoted theorem, proof sheet, finite semantic case, replay table, comparator row, or claim-boundary section, the source remains background support rather than public theorem evidence.


For Processes: Processes K8, the audit route keeps the full file `content/processes/processes\_k8.tex` indexed by the science monolith corpus register and the public evidence package. The science is therefore not discarded, while the PDF avoids becoming a raw source dump. A line-level reviewer follows the register; a scientific reader follows this digest and the main chapter.


Processes: Processes K8 reader checkpoint is satisfied only when the main manuscript tells the reader which claim class this module can support, which evidence class limits it, which prior-art or comparator boundary applies, and which failure would reopen the claim. This fourth check is deliberately positive: the module must contribute intelligibility, not merely avoid forbidden wording.


For Processes: Processes K8, the external wording boundary allows public language to say that the route helps organize, test, or constrain a bounded OC claim. It may not say that the route proves full-domain closure, final full-scope synthesis status, scoped unrestricted comparative claim, or a domain result absent from the current proof and evidence layer. This boundary is part of the source's scientific content, not a marketing caveat.


\paragraph{Processes: Processes K9}

For Processes: Processes K9, the scholarly synthesis contributes process schema, ordering, and change conditions, explains why this module exists, shows how it connects to the current 1.3.3 argument, and prevents stronger inherited language from being promoted without the current proof, replay, comparator, and falsifier surfaces.


For Processes: Processes K9, reader use is to ask whether the process claim names the state, transition, boundary, and endpoint condition that make it testable; if that source-specific question cannot be answered from the promoted theorem, proof sheet, finite semantic case, replay table, comparator row, or claim-boundary section, the source remains background support rather than public theorem evidence.


For Processes: Processes K9, the audit route keeps the full file `content/processes/processes\_k9.tex` indexed by the science monolith corpus register and the public evidence package. The science is therefore not discarded, while the PDF avoids becoming a raw source dump. A line-level reviewer follows the register; a scientific reader follows this digest and the main chapter.


Processes: Processes K9 reader checkpoint is satisfied only when the main manuscript tells the reader which claim class this module can support, which evidence class limits it, which prior-art or comparator boundary applies, and which failure would reopen the claim. This fourth check is deliberately positive: the module must contribute intelligibility, not merely avoid forbidden wording.


For Processes: Processes K9, the external wording boundary allows public language to say that the route helps organize, test, or constrain a bounded OC claim. It may not say that the route proves full-domain closure, final full-scope synthesis status, scoped unrestricted comparative claim, or a domain result absent from the current proof and evidence layer. This boundary is part of the source's scientific content, not a marketing caveat.


\paragraph{Processes: Processes K10}

For Processes: Processes K10, the scholarly synthesis contributes process schema, ordering, and change conditions, explains why this module exists, shows how it connects to the current 1.3.3 argument, and prevents stronger inherited language from being promoted without the current proof, replay, comparator, and falsifier surfaces.


For Processes: Processes K10, reader use is to ask whether the process claim names the state, transition, boundary, and endpoint condition that make it testable; if that source-specific question cannot be answered from the promoted theorem, proof sheet, finite semantic case, replay table, comparator row, or claim-boundary section, the source remains background support rather than public theorem evidence.


For Processes: Processes K10, the audit route keeps the full file `content/processes/processes\_k10.tex` indexed by the science monolith corpus register and the public evidence package. The science is therefore not discarded, while the PDF avoids becoming a raw source dump. A line-level reviewer follows the register; a scientific reader follows this digest and the main chapter.


Processes: Processes K10 reader checkpoint is satisfied only when the main manuscript tells the reader which claim class this module can support, which evidence class limits it, which prior-art or comparator boundary applies, and which failure would reopen the claim. This fourth check is deliberately positive: the module must contribute intelligibility, not merely avoid forbidden wording.


For Processes: Processes K10, the external wording boundary allows public language to say that the route helps organize, test, or constrain a bounded OC claim. It may not say that the route proves full-domain closure, final full-scope synthesis status, scoped unrestricted comparative claim, or a domain result absent from the current proof and evidence layer. This boundary is part of the source's scientific content, not a marketing caveat.


\paragraph{Processes: Processes K11}

For Processes: Processes K11, the scholarly synthesis contributes process schema, ordering, and change conditions, explains why this module exists, shows how it connects to the current 1.3.3 argument, and prevents stronger inherited language from being promoted without the current proof, replay, comparator, and falsifier surfaces.


For Processes: Processes K11, reader use is to ask whether the process claim names the state, transition, boundary, and endpoint condition that make it testable; if that source-specific question cannot be answered from the promoted theorem, proof sheet, finite semantic case, replay table, comparator row, or claim-boundary section, the source remains background support rather than public theorem evidence.


For Processes: Processes K11, the audit route keeps the full file `content/processes/processes\_k11.tex` indexed by the science monolith corpus register and the public evidence package. The science is therefore not discarded, while the PDF avoids becoming a raw source dump. A line-level reviewer follows the register; a scientific reader follows this digest and the main chapter.


Processes: Processes K11 reader checkpoint is satisfied only when the main manuscript tells the reader which claim class this module can support, which evidence class limits it, which prior-art or comparator boundary applies, and which failure would reopen the claim. This fourth check is deliberately positive: the module must contribute intelligibility, not merely avoid forbidden wording.


For Processes: Processes K11, the external wording boundary allows public language to say that the route helps organize, test, or constrain a bounded OC claim. It may not say that the route proves full-domain closure, final full-scope synthesis status, scoped unrestricted comparative claim, or a domain result absent from the current proof and evidence layer. This boundary is part of the source's scientific content, not a marketing caveat.


\paragraph{Processes: Processes K12}

For Processes: Processes K12, the scholarly synthesis contributes process schema, ordering, and change conditions, explains why this module exists, shows how it connects to the current 1.3.3 argument, and prevents stronger inherited language from being promoted without the current proof, replay, comparator, and falsifier surfaces.


For Processes: Processes K12, reader use is to ask whether the process claim names the state, transition, boundary, and endpoint condition that make it testable; if that source-specific question cannot be answered from the promoted theorem, proof sheet, finite semantic case, replay table, comparator row, or claim-boundary section, the source remains background support rather than public theorem evidence.


For Processes: Processes K12, the audit route keeps the full file `content/processes/processes\_k12.tex` indexed by the science monolith corpus register and the public evidence package. The science is therefore not discarded, while the PDF avoids becoming a raw source dump. A line-level reviewer follows the register; a scientific reader follows this digest and the main chapter.


Processes: Processes K12 reader checkpoint is satisfied only when the main manuscript tells the reader which claim class this module can support, which evidence class limits it, which prior-art or comparator boundary applies, and which failure would reopen the claim. This fourth check is deliberately positive: the module must contribute intelligibility, not merely avoid forbidden wording.


For Processes: Processes K12, the external wording boundary allows public language to say that the route helps organize, test, or constrain a bounded OC claim. It may not say that the route proves full-domain closure, final full-scope synthesis status, scoped unrestricted comparative claim, or a domain result absent from the current proof and evidence layer. This boundary is part of the source's scientific content, not a marketing caveat.


\subsubsection{Content}

\paragraph{Content: Operators scoped}

For Content: Operators scoped, the scholarly synthesis contributes operator semantics and update admissibility, explains why this module exists, shows how it connects to the current 1.3.3 argument, and prevents stronger inherited language from being promoted without the current proof, replay, comparator, and falsifier surfaces.


For Content: Operators scoped, reader use is to ask whether the operator is a smooth flow, a guard/reset update, a rewrite, or a typed non-smooth transformation; if that source-specific question cannot be answered from the promoted theorem, proof sheet, finite semantic case, replay table, comparator row, or claim-boundary section, the source remains background support rather than public theorem evidence.


For Content: Operators scoped, the audit route keeps the full file `content/operators\_universal.tex` indexed by the science monolith corpus register and the public evidence package. The science is therefore not discarded, while the PDF avoids becoming a raw source dump. A line-level reviewer follows the register; a scientific reader follows this digest and the main chapter.


Content: Operators scoped reader checkpoint is satisfied only when the main manuscript tells the reader which claim class this module can support, which evidence class limits it, which prior-art or comparator boundary applies, and which failure would reopen the claim. This fourth check is deliberately positive: the module must contribute intelligibility, not merely avoid forbidden wording.


For Content: Operators scoped, the external wording boundary allows public language to say that the route helps organize, test, or constrain a bounded OC claim. It may not say that the route proves full-domain closure, final full-scope synthesis status, scoped unrestricted comparative claim, or a domain result absent from the current proof and evidence layer. This boundary is part of the source's scientific content, not a marketing caveat.


Synthesis checkpoint after 84 indexed modules in content. The preceding modules teach inspection logic; they do not inflate claim strength. The release remains bounded by the current proof and evidence surface.


\paragraph{Content: Theorems Master}

For Content: Theorems Master, the scholarly synthesis contributes historical theorem route and theorem-boundary pressure, explains why this module exists, shows how it connects to the current 1.3.3 argument, and prevents stronger inherited language from being promoted without the current proof, replay, comparator, and falsifier surfaces.


For Content: Theorems Master, reader use is to ask whether the theorem label is current proof, historical route, definition, or future proof obligation; if that source-specific question cannot be answered from the promoted theorem, proof sheet, finite semantic case, replay table, comparator row, or claim-boundary section, the source remains background support rather than public theorem evidence.


For Content: Theorems Master, the audit route keeps the full file `content/theorems\_master.tex` indexed by the science monolith corpus register and the public evidence package. The science is therefore not discarded, while the PDF avoids becoming a raw source dump. A line-level reviewer follows the register; a scientific reader follows this digest and the main chapter.


Content: Theorems Master reader checkpoint is satisfied only when the main manuscript tells the reader which claim class this module can support, which evidence class limits it, which prior-art or comparator boundary applies, and which failure would reopen the claim. This fourth check is deliberately positive: the module must contribute intelligibility, not merely avoid forbidden wording.


For Content: Theorems Master, the external wording boundary allows public language to say that the route helps organize, test, or constrain a bounded OC claim. It may not say that the route proves full-domain closure, final full-scope synthesis status, scoped unrestricted comparative claim, or a domain result absent from the current proof and evidence layer. This boundary is part of the source's scientific content, not a marketing caveat.


\paragraph{Content: Complexity S}

For Content: Complexity S, the scholarly synthesis contributes complexity-functional vocabulary and limits, explains why this module exists, shows how it connects to the current 1.3.3 argument, and prevents stronger inherited language from being promoted without the current proof, replay, comparator, and falsifier surfaces.


For Content: Complexity S, reader use is to ask whether a complexity statement names the functional and transition before claiming increase or comparison; if that source-specific question cannot be answered from the promoted theorem, proof sheet, finite semantic case, replay table, comparator row, or claim-boundary section, the source remains background support rather than public theorem evidence.


For Content: Complexity S, the audit route keeps the full file `content/complexity\_S.tex` indexed by the science monolith corpus register and the public evidence package. The science is therefore not discarded, while the PDF avoids becoming a raw source dump. A line-level reviewer follows the register; a scientific reader follows this digest and the main chapter.


Content: Complexity S reader checkpoint is satisfied only when the main manuscript tells the reader which claim class this module can support, which evidence class limits it, which prior-art or comparator boundary applies, and which failure would reopen the claim. This fourth check is deliberately positive: the module must contribute intelligibility, not merely avoid forbidden wording.


For Content: Complexity S, the external wording boundary allows public language to say that the route helps organize, test, or constrain a bounded OC claim. It may not say that the route proves full-domain closure, final full-scope synthesis status, scoped unrestricted comparative claim, or a domain result absent from the current proof and evidence layer. This boundary is part of the source's scientific content, not a marketing caveat.
